\documentclass[11pt,a4paper]{article}
\usepackage[utf8]{inputenc}
\usepackage[T1]{fontenc}
\usepackage[czech]{babel}
\usepackage[margin=2.0cm,top=2.0cm,bottom=2.2cm]{geometry}
\usepackage{amsmath,amssymb,amsthm}
\usepackage{parskip}
\usepackage{enumitem}
\setlist{nosep,leftmargin=1.4em}
\usepackage{booktabs}
\usepackage{array}
\usepackage{xcolor}
\definecolor{linkblue}{RGB}{20,60,150}
\definecolor{fmlcol}{RGB}{20,70,140}
\definecolor{defbg}{RGB}{240,244,252}
\definecolor{defframe}{RGB}{100,130,190}
\definecolor{markcol}{RGB}{176,84,0}
\definecolor{markbg}{RGB}{253,246,236}
\usepackage[most]{tcolorbox}
\newtcolorbox{poznbox}[1][]{breakable,colback=markbg,colframe=markcol,boxrule=0.6pt,arc=1.5mm,left=2mm,right=2mm,top=1mm,bottom=1mm,title={#1},fonttitle=\bfseries}
\usepackage[colorlinks=true,linkcolor=linkblue,urlcolor=linkblue]{hyperref}

% --- sazba vzorců: název (\fml) + display math + popis (\pop) ---
\setlength{\abovedisplayskip}{4pt plus 2pt minus 1pt}
\setlength{\belowdisplayskip}{4pt plus 2pt minus 1pt}
\setlength{\abovedisplayshortskip}{1pt plus 1pt}
\setlength{\belowdisplayshortskip}{4pt plus 1pt}
\setlength{\parskip}{0.35em}

\newcommand{\fml}[1]{\par\addvspace{2.6mm}\noindent
  \textcolor{fmlcol}{$\blacktriangleright$}\ \textbf{#1}\par\nobreak}
\newcommand{\pop}[1]{\par\nobreak\noindent{\small #1}\par}

% oddělovač okruhů: #1 = číslo zdrojového PDF, #2 = název okruhu, #3 = podtitul
\newcommand{\okruh}[3]{%
  \clearpage
  \phantomsection
  \addcontentsline{toc}{part}{#1\quad #2}%
  \begingroup\centering
    {\Huge\bfseries #1\quad #2\par}\vspace{1.5mm}%
    {\small\itshape #3\par}%
  \endgroup
  \vspace{4mm}}

% Kapitoly a~podkapitoly nesou ČÍSLO ZE ZDROJOVÉHO DOKUMENTU (01/02/03), nikoli
% průběžné. Mezery v~číslování znamenají, že daná kapitola žádný vzorec nemá.
\newcommand{\kap}[2]{\setcounter{section}{#1}\addtocounter{section}{-1}\section{#2}}
\newcommand{\podkap}[2]{\setcounter{subsection}{#1}\addtocounter{subsection}{-1}\subsection{#2}}
\newcommand{\podpodkap}[2]{\setcounter{subsubsection}{#1}\addtocounter{subsubsection}{-1}\subsubsection{#2}}

% Nadpisů je hodně (kopírují členění zdrojů), proto sevřenější vertikální mezery.
\makeatletter
\renewcommand\section{\@startsection{section}{1}{\z@}%
  {-2.4ex \@plus -1ex \@minus -.2ex}{1.0ex \@plus.2ex}{\normalfont\Large\bfseries}}
\renewcommand\subsection{\@startsection{subsection}{2}{\z@}%
  {-2.0ex \@plus -1ex \@minus -.2ex}{0.6ex \@plus.2ex}{\normalfont\large\bfseries}}
\renewcommand\subsubsection{\@startsection{subsubsection}{3}{\z@}%
  {-1.7ex \@plus -1ex \@minus -.2ex}{0.4ex \@plus.2ex}{\normalfont\normalsize\bfseries}}
\makeatother

\setcounter{tocdepth}{2}
\setcounter{secnumdepth}{3}

\title{\textbf{Vzorce a~rovnice ke~SZZ}\\[2mm]
\large Multiagentní systémy \textbullet{} Přírodou inspirované počítání \textbullet{} Dobývání znalostí}
\author{SZZ Umělá inteligence, MFF UK}
\date{srpen 2026}

\begin{document}
\maketitle

\begin{center}
\begin{minipage}{0.93\textwidth}
\small
Výcuc ze tří učebních textů (\texttt{01-multiagentni-systemy}, \texttt{02-prirodou-in\-spi\-ro\-va\-ne-pocitani},
\texttt{03-dobyvani-znalosti}): \textbf{jen vzorce a~rovnice, které je potřeba umět}, každý
s~krátkým popisem toho, co znamená a~kde se používá. Slouží k~poslednímu opakování
a~k~ověření, že člověk formuli umí napsat bez nahlédnutí do textu --- \emph{proč} to tak
je a~kdy se to používá je v~příslušné kapitole zdrojového textu.
\end{minipage}
\end{center}

\begin{center}
\begin{minipage}{0.93\textwidth}
\begin{poznbox}[Jak text používat]
\small
\textbf{Názvy i~čísla kapitol a~podkapitol jsou převzaty ze zdrojových textů}
\texttt{01}, \texttt{02} a~\texttt{03}, takže se dá listovat oběma směry: kapitola
\emph{3.2.4 Support Vector Machines} je zde i~tam pod stejným číslem. \emph{Mezery
v~číslování} znamenají, že daná kapitola zdroje žádný vzorec neobsahuje. Vzorec, který
pokrývá dvě sousední podkapitoly zdroje, je zařazen pod tu první a~pojmenován podle obou.
U~každého vzorce je uveden význam všech symbolů; kde je vzorců k~jednomu tématu více, jsou
v~pořadí, v~jakém se odvozují nebo používají.

\textbf{Značení sdílené napříč celým textem:} $f$ fitness / hustota,
tříd/atributů, $P$ populace nebo pravděpodobnost, $\gamma$ diskontní faktor,
$\alpha$ rychlost učení nebo hladina významnosti, $\eta$ rychlost učení.
Význam se vždy upřesňuje na místě.
\end{poznbox}
\end{minipage}
\end{center}

\tableofcontents

% ==========================================================================
%% ========================================================================
\okruh{01}{Multiagentní systémy}{Abstraktní architektury, plánování, hledání cesty, teorie her a~aukce, zpětnovazební učení.}
%% ========================================================================

\kap{1}{Architektura autonomního agenta a~mechanismus výběru akcí}

\podkap{5}{Abstraktní architektura agenta}

\fml{Běh agenta v~prostředí}
\[
r \;=\; e_0 \xrightarrow{a_0} e_1 \xrightarrow{a_1} e_2 \xrightarrow{a_2} \cdots
\xrightarrow{a_{n-1}} e_n
\]
\pop{$E = \{e, e', \dots\}$ je konečná množina stavů prostředí, $A = \{a, a', \dots\}$
konečná množina akcí agenta. \emph{Běh} je posloupnost střídajících se stavů a~akcí.
$R$ je množina všech konečných běhů, $R^A \subseteq R$ ty končící akcí,
$R^E \subseteq R$ ty končící stavem prostředí.}

\fml{Přechodová funkce prostředí (\emph{state transformer function})}
\[
\tau : R^A \to 2^E, \qquad \mathit{Env} = \langle E, e_0, \tau\rangle
\]
\pop{Zobrazuje běh končící akcí na \emph{množinu} možných následujících stavů --- prostředí je
tedy nedeterministické a~\emph{historicky závislé}. $\tau(r) = \emptyset$ znamená konec běhu.
Prostředí je trojice: stavy, počáteční stav, přechodová funkce.}

\fml{Agent jako funkce běhu}
\[
Ag : R^E \to A
\]
\pop{Agent rozhoduje deterministicky a~na základě \emph{celé historie} systému;
$\mathcal{AG}$ značí množinu všech agentů.}

\podkap{6}{Čistě reaktivní agent a~agent s~vnitřním stavem}

\fml{Čistě reaktivní agent (\emph{tropistic agent})}
\[
Ag : E \to A
\]
\pop{Reaguje pouze na \emph{aktuální} stav prostředí, historii ignoruje. Ke~každému
reaktivnímu agentovi existuje funkčně ekvivalentní standardní agent, \emph{obráceně to
neplatí}. Termostat: $Ag(e) = \mathrm{off}$ pro \uv{teplota OK}, jinak $\mathrm{on}$.}

\fml{Agent s~vnitřním stavem}
\begin{align*}
\mathit{see} &: E \to \mathit{Per} && \text{(vnímání)}\\
\mathit{next} &: I \times \mathit{Per} \to I && \text{(aktualizace vnitřního stavu)}\\
\mathit{action} &: I \to A && \text{(výběr akce)}
\end{align*}
\pop{$I$ je množina vnitřních stavů, $\mathit{Per}$ množina vjemů. Cyklus: agent startuje
v~$i_0$, \emph{see} převede stav prostředí na vjem, \emph{next} aktualizuje vnitřní stav,
\emph{action} zvolí akci. Vyjadřovací síla je \emph{stejná} jako u~standardního agenta.}

\podkap{7}{Utilita: jak agentovi říct, co má dělat}

\fml{Utilita stavu a~utilita běhu}
\[
u : E \to \mathbb{R}, \qquad\qquad u : R \to \mathbb{R}
\]
\pop{Cíl se agentovi zadává \emph{nepřímo}, mírou úspěšnosti. Hodnocení jednotlivých stavů je
krátkozraké, proto se zavádí utilita celého běhu (např.\ utilita nejhoršího navštíveného stavu
nebo průměr).}

\fml{Optimální agent (maximalizace očekávané utility)}
\[
Ag_{\mathrm{opt}} \;=\; \arg\max_{Ag \in \mathcal{AG}} \sum_{r \in R(Ag,\mathit{Env})}
u(r)\, P(r \mid Ag, \mathit{Env}),
\qquad \sum_{r} P(r \mid Ag,\mathit{Env}) = 1
\]
\pop{$P(r \mid Ag, \mathit{Env})$ je pravděpodobnost běhu $r$. Definice je elegantní, ale
\textbf{nekonstruktivní} --- neříká, jak takového agenta navrhnout, a~utilitní funkci bývá
obtížné vůbec zadat.}

\fml{Tileworld --- utilita běhu}
\[
u(r) = \frac{N_f}{N_a}
\]
\pop{$N_f$ = počet děr zaplněných agentem během běhu $r$, $N_a$ = počet všech děr, které se
během $r$ objevily. Klasický příklad utility v~\emph{dynamickém} prostředí.}

\kap{2}{Metody pro řízení agentů: symbolické a~konekcionistické reaktivní plánování, hybridní přístupy}

\podkap{1}{Klasický přístup UI a~deduktivní agent}

\fml{Deduktivní agent (agent jako dokazovač vět)}
\[
\mathit{DB} \vdash_\rho \varphi
\]
\pop{$L$ je množina formulí predikátové logiky 1.~řádu, $D = 2^L$ množina databází,
vnitřní stav agenta je $\mathit{DB} \in D$. Zápis znamená \uv{formuli $\varphi$ lze odvodit
z~$\mathit{DB}$ pomocí pravidel $\rho$}. Agent hledá akci, pro niž lze odvodit
$\mathit{Do}(a)$; nejde-li to, akci alespoň \emph{konzistentní} s~databází.}

\podkap{3}{Plánování: STRIPS}

\fml{STRIPS --- posloupnost databází podle plánu}
\[
B_i = (B_{i-1} \setminus D_{a_i}) \cup A_{a_i}, \qquad i = 1,\dots,n
\]
\pop{Akce je popsána trojicí $\langle P_a, D_a, A_a\rangle$ (předpoklady, mazaná množina,
přidávaná množina); plánovací problém je $\langle B_0, O, G\rangle$.
Plán $\pi = (a_1,\dots,a_n)$ je \textbf{přípustný}, pokud $B_{i-1} \models P_{a_i}$ pro
každé $i$, a~\textbf{korektní}, je-li přípustný a~navíc $B_n \models G$.}

\kap{3}{Problém hledání cesty}

\podkap{2}{A*}

\fml{Ohodnocení vrcholu v~A*}
\begin{align*}
g(n) &= \text{cena nejlepší dosud nalezené cesty } s \to n, \\
h(n) &= \text{heuristický odhad ceny } n \to g_{\mathrm{cil}}, \\
f(n) &= g(n) + h(n)
\end{align*}
\pop{$f(n)$ je odhad ceny nejlepší cesty $s \to g_{\mathrm{cil}}$ \emph{vedoucí přes} $n$;
A* expanduje vrchol s~nejmenším $f$. Pro $h \equiv 0$ degeneruje na Dijkstru.}

\fml{Přípustnost a~konzistence heuristiky}
\[
h(n) \le h^*(n) \quad\text{(přípustná)}, \qquad
h(n) \le w(n,m) + h(m),\ \ h(g_{\mathrm{cil}}) = 0 \quad\text{(konzistentní)}
\]
\pop{$h^*(n)$ je skutečná cena nejlevnější cesty z~$n$ do cíle; přípustná heuristika
\emph{nikdy nenadhodnocuje}. Konzistence je trojúhelníková nerovnost na každé hraně $(n,m)$ s~cenou $w(n,m)$.
Platí: \textbf{konzistence $\Rightarrow$ přípustnost}. Při přípustné $h$ je A* optimální,
při konzistentní $h$ je navíc $f$ podél cesty neklesající (stačí uzavřený seznam bez
reexpanzí).}

\podkap{3}{Dynamické prostředí: D*~Lite}

\fml{LPA* / D*~Lite --- lookahead hodnota}
\[
\mathit{rhs}(n) = \begin{cases}
0 & n = s,\\
\min_{p \in \mathit{pred}(n)} \big(g(p) + w(p,n)\big) & \text{jinak}
\end{cases}
\]
\pop{$\mathit{pred}(n)$ jsou předchůdci $n$, $w(p,n)$ cena hrany $(p,n)$. Vrchol je \emph{lokálně konzistentní}, pokud $g(n) = \mathit{rhs}(n)$. Změní-li se cena
hrany, stanou se některé vrcholy nekonzistentními, zařadí se do prioritní fronty a~změna se
propaguje jen do zasažené části --- proto inkrementální přeplánování stojí méně než A*
od~nuly.}

\kap{4}{Komunikace a~znalosti v~MAS: ontologie, řečové akty, FIPA-ACL, protokoly}

\podkap{7}{KQML a~FIPA-ACL}

\fml{Semantika performativu \emph{inform} (FP / RE)}
\[
\text{FP: } B_i\varphi \wedge \neg B_i\big(\mathit{Bif}_j\varphi \vee \mathit{Uif}_j\varphi\big),
\qquad
\text{RE: } B_j\varphi
\]
\pop{Sémantika ACL se zapisuje v~modální logice s~operátory $B_i\varphi$ (agent $i$ věří
$\varphi$), $U_i\varphi$ (je nejistý), $I_i\varphi$ (má záměr).
\textbf{FP} (\emph{feasibility precondition}) --- co musí platit, aby odesílatel mohl akt
provést; \textbf{RE} (\emph{rational effect}) --- efekt, kvůli němuž akt provádí.
Příjemce \emph{není povinen} RE naplnit (autonomie!).}

\kap{5}{Distribuované řešení problémů, kooperace, teorie her a~aukce}

\podkap{8}{Hra v~normální formě}

\fml{Hra v~normální formě}
\[
\langle N, (A_i)_{i \in N}, (u_i)_{i \in N}\rangle
\]
\pop{$N$ množina hráčů, $A_i$ množina akcí hráče $i$, $u_i$ jeho utilitní funkce.
Pro dva hráče se zapisuje \emph{výplatní maticí} (řádky = strategie $i$, sloupce = strategie
$j$, buňka $u_i / u_j$). Strategie je \emph{ryzí} (jedna akce) nebo \emph{smíšená}
(rozdělení pravděpodobnosti nad ryzími).}

\fml{Nashovo ekvilibrium}
\[
u_i(s_i^*, s_{-i}^*) \;\ge\; u_i(s_i, s_{-i}^*) \qquad \text{pro každého } i
\text{ a~každou jeho strategii } s_i
\]
\pop{$s_{-i}$ značí strategie všech ostatních hráčů. Profil $s^*$ je NE, pokud
\textbf{žádný hráč nemá jednostranný důvod se odchýlit} --- strategie jsou vzájemně nejlepší
odpovědi. V~ryzích strategiích nemusí existovat, ve smíšených existuje vždy (Nash).}

\podkap{9}{Klasické hry}

\fml{Vězňovo dilema --- podmínky na výplaty}
\[
T > R > P > S \qquad\text{a často navíc}\qquad 2R > T + S
\]
\pop{$R$ (\emph{reward}) za oboustrannou spolupráci, $P$ (\emph{punishment}) za oboustrannou
zradu, $T$ (\emph{temptation}) za zradu spolupracujícího, $S$ (\emph{sucker's payoff})
za~zrazení. $T>R$ a~$P>S$ činí zradu \emph{dominantní strategií}, takže $DD$ je jediné NE
--- ale $CC$ je Pareto-lepší. Druhá podmínka ($2R > T+S$) je nutná pro iterovanou verzi:
střídavé vykořisťování se nesmí vyplatit více než trvalá spolupráce.}

\fml{Výpočet smíšeného NE ve hře $2\times2$ (podmínka indiference)}
\[
u_i(\text{akce}_1) = u_i(\text{akce}_2)
\qquad\text{např.}\qquad
1\cdot q + 0\cdot(1-q) \;=\; 0\cdot q + 2\cdot(1-q)
\]
\pop{Protihráč $j$ volí pravděpodobnost $q$ tak, aby byl hráč $i$ \emph{indiferentní} mezi
svými akcemi (jinak by $i$ hrál ryzí strategii). Z~uvedené rovnosti $q = 2-2q$, tedy
$q = 2/3$. V~\uv{Bitvě pohlaví} tak každý hraje svou preferovanou možnost s~$2/3$
a~očekávaná utilita $2/3$ je \emph{nižší} než v~kterémkoli ryzím ekvilibriu.}

\podkap{15}{Návrh mechanismů, VCG a~vyjednávání}

\fml{VCG --- Vickrey--Clarke--Groves}
\[
p_i \;=\; \max_{x} \sum_{j \ne i} \hat v_j(x) \;-\; \sum_{j \ne i} \hat v_j(x^*)
\]
\pop{$\hat v_i$ je ohlášená valuace, $x^*$ alokace maximalizující $\sum_i \hat v_i(x)$.
Agent $i$ platí \textbf{externalitu, kterou způsobuje ostatním} --- rozdíl mezi tím, co by
ostatní získali bez něj, a~tím, co získají s~ním. Mechanismus je pravdomluvný, efektivní
a~individuálně racionální; pro jeden předmět se redukuje na Vickreyovu aukci
(2.~cena). Nevýhody: WDP se řeší $n{+}1$-krát, výnos prodávajícího může být i~nulový.}

\fml{Zeuthenova strategie --- ochota riskovat konflikt}
\[
\mathrm{Risk}_i \;=\; \frac{u_i(\text{vlastní návrh}) - u_i(\text{návrh protějšku})}
{u_i(\text{vlastní návrh})}
\]
\pop{V~monotónním protokolu ústupků ustupuje agent s~\emph{nižším} $\mathrm{Risk}$, a~to
právě tak málo, aby se poměry obrátily. Čitatel je \uv{co ztratím přijetím cizího návrhu},
jmenovatel \uv{co mám v~sázce}. Dvojice Zeuthenových strategií je v~Nashově ekvilibriu
a~dohoda maximalizuje součin utilit (Nashovo řešení).}

\kap{6}{Metody pro učení agentů a~zpětnovazební učení}

\podkap{2}{Markovský rozhodovací proces}

\fml{Cíl v~MDP --- očekávaný diskontovaný výnos}
\[
\mathbb{E}\Big[ \sum_{t=0}^{\infty} \gamma^t R(s_t, a_t, s_{t+1}) \Big]
\]
\pop{MDP je $\langle S, A, T, R\rangle$ s~přechodovými pravděpodobnostmi $T(s,a,s')$
a~odměnou $R$; platí \emph{markovská vlastnost} (další stav a~odměna závisí jen na aktuálním
stavu a~akci). $\gamma \in [0,1)$ diskontuje budoucnost a~zaručuje konečnost součtu.
Politika je $\pi : S \to A$ (nebo stochasticky $\pi(a|s)$); hledáme $\pi^*$ maximalizující
tento výnos.}

\fml{Bellmanova rovnice pro evaluaci politiky}
\[
V^\pi(s) = \sum_{s'} T(s,\pi(s),s') \big[ R(s,\pi(s),s') + \gamma V^\pi(s') \big]
\]
\pop{$V^\pi(s)$ je očekávaná utilita stavu $s$, sleduje-li agent politiku $\pi$;
$Q^\pi(s,a)$ totéž pro dvojici stav--akce. Rovnice je rekurzivní rozklad
\uv{okamžitá odměna $+$ diskontovaná hodnota následníka}.}

\fml{Bellmanovy rovnice optimality}
\begin{align*}
V^*(s) &= \max_{a \in A} \sum_{s'} T(s,a,s') \big[ R(s,a,s') + \gamma V^*(s')\big], \\
Q^*(s,a) &= \sum_{s'} T(s,a,s')\big[R(s,a,s') + \gamma \max_{a'} Q^*(s',a')\big]
\end{align*}
\pop{Optimální politika se získá jako $\pi^*(s) = \arg\max_a Q^*(s,a)$ --- z~$Q^*$ tedy
\emph{bez} znalosti modelu, z~$V^*$ jen se znalostí $T$ a~$R$. Iterováním těchto rovnic
pracují \emph{value iteration} a~\emph{policy iteration} (model-based).}

\podkap{3}{Q-learning a~SARSA}

\fml{Q-learning (Watkins, 1989)}
\[
Q(s,a) \;\leftarrow\; Q(s,a) + \alpha\Big[\, \underbrace{r + \gamma \max_{a'} Q(s',a')}
_{\text{TD cíl}} - Q(s,a) \Big],
\qquad 0 \le \alpha \le 1
\]
\pop{Agent provedl ve stavu $s$ akci $a$, dostal \textbf{okamžitou odměnu} $r$
(jedno pozorování, nikoli funkce $R$ jako u~Bellmana) a~přešel do $s'$; $a'$ probíhá akce
dostupné v~$s'$, $\alpha$ je rychlost učení a~$\gamma$ diskontní faktor. Ekvivalentně $Q(s,a) \leftarrow (1-\alpha)Q(s,a) + \alpha\big(r + \gamma\max_{a'}Q(s',a')\big)$;
hranatá závorka je \emph{TD chyba}.
\textbf{Off-policy}: cíl používá $\max_{a'}$, učí se tedy hodnotu \emph{optimální} politiky
bez ohledu na to, jak agent skutečně jedná. Konverguje $Q \to Q^*$, je-li každá dvojice
$(s,a)$ navštívena nekonečněkrát a~$\alpha$ vhodně klesá.}

\fml{SARSA}
\[
Q(s,a) \leftarrow Q(s,a) + \alpha\big[r + \gamma\, Q(s',a') - Q(s,a)\big]
\]
\pop{Aktualizace používá \emph{skutečně zvolenou} následující akci $a'$ (podle stejné
politiky, kterou agent jedná) --- \textbf{on-policy}. Učí se hodnotu politiky včetně
náhodného průzkumu, a~je proto \emph{opatrnější}: v~úloze \uv{chůze po útesu} se drží dál
od~propasti než Q-learning.}

\fml{$\varepsilon$-greedy a~softmax (dilema průzkumu a~využití)}
\[
\pi(s) = \begin{cases}
\text{náhodná akce z } A & \text{s pravděpodobností } \varepsilon,\\[2pt]
\arg\max_a Q(s,a) & \text{s pravděpodobností } 1-\varepsilon,
\end{cases}
\qquad
\pi(a|s) \propto e^{Q(s,a)/\tau}
\]
\pop{$\varepsilon$ se v~čase snižuje (\emph{annealing}). Softmax/Boltzmann s~teplotou $\tau$
váží akce podle jejich hodnoty (ne uniformně jako $\varepsilon$-greedy).
Další alternativy: optimistická inicializace $Q$, UCB.}

\podkap{4}{Metody založené na politice}

\fml{REINFORCE (policy gradient)}
\[
z_{t+1} = z_t + \alpha\, G_t\, \nabla_z \ln \pi(A_t \mid S_t, z_t)
\]
\pop{Politika je parametrizovaná vektorem $z$, tedy $\pi(a \mid s, z)$; $G_t$ je výnos od času
$t$ (Monte Carlo z~celé epizody). Intuice: \emph{zvyš pravděpodobnost akcí, po nichž
následoval vysoký výnos}. Zvládá \emph{spojité} akční prostory a~stochastické politiky, ale
má velký rozptyl (řeší se odečtením baseline --- \emph{actor--critic}).}

%% ========================================================================
\okruh{02}{Přírodou inspirované počítání}{Selekce a~operátory, teorie schémat, evoluční strategie, diferenciální evoluce, rojové algoritmy, lokální prohledávání, aplikace.}
%% ========================================================================

\kap{1}{Genetické algoritmy, genetické a~evoluční programování}

\podkap{1}{Obecné schéma evolučního algoritmu}

\podpodkap{5}{Selekční mechanismy}

\fml{Ruletová (fitness-proporcionální) selekce}
\[
p_i \;=\; \frac{f(x_i)}{\sum_{j=1}^{n} f(x_j)} \;=\; \frac{f(x_i)}{n\,\bar{f}},
\qquad
\mathbb{E}[\text{počet výběrů } x_i] = n\,p_i = \frac{f(x_i)}{\bar f}
\]
\pop{$x_i$ je $i$-tý jedinec, $\bar f$ průměrná fitness populace, $n$ její velikost; výběr je \emph{s~vracením},
takže jeden jedinec může být vybrán vícekrát. Nevýhody: předčasná dominance
\uv{supermana} (jedinec s~extrémní fitness ovládne populaci a~diverzita zmizí), ztráta
selekčního tlaku na konci běhu (podobné fitness $\Rightarrow$ téměř náhodný výběr),
citlivost na afinní transformace fitness ($f$ a~$f+c$ dají jiné $p_i$) a~vysoký rozptyl
skutečných počtů výběrů.}

\fml{Turnajová selekce velikosti $k$}
\[
p_i \;=\; \frac{(n-i+1)^k - (n-i)^k}{n^k}
\qquad (i = 1 \text{ značí \textbf{nejlepšího}})
\]
\pop{Vyber náhodně (uniformně, s~vracením) $k$ jedinců a~vrať nejlepšího; $k$ (typicky 2,~3,~5)
přímo řídí selekční tlak. Odvození: pravděpodobnost, že všech $k$ padne mezi $i$-tého
a~horší, mínus pravděpodobnost, že všech $k$ padne ostře pod $i$-tého.
Výhody: nepotřebuje globální informaci, je invariantní vůči monotónním transformacím fitness,
paralelizovatelná, funguje i~tam, kde fitness není dána a~porovnání se získává souhrou.}

\fml{Lineární pořadová selekce}
\[
p_i \;=\; \frac{1}{n}\Big(s^{-} + (s^{+}-s^{-})\,\frac{i-1}{n-1}\Big),
\qquad 1 \le s^{+} \le 2,\quad s^{-} = 2 - s^{+}
\]
\pop{\emph{Pozor na konvenci:} zde $i=1$ značí \textbf{nejhoršího}, $i=n$ nejlepšího.
$s^{+}$ je očekávaný počet kopií nejlepšího jedince. Exponenciální varianta:
$p_i \propto (1-c)\,c^{\,n-i}$. Pořadová selekce udržuje \emph{konstantní} selekční tlak
během celého běhu a~je odolná vůči \uv{supermanům} i~vůči změně měřítka fitness.}

\fml{Boltzmannova selekce}
\[
p_i \;\propto\; e^{f(x_i)/T}
\]
\pop{Teplota $T$ klesá v~čase: vysoké $T$ $\Rightarrow$ téměř uniformní selekce (explorace),
nízké $T$ $\Rightarrow$ téměř deterministická (exploatace). Přímá souvislost se simulovaným
žíháním.}

\fml{Doba převzetí (\emph{takeover time})}
\[
\tau^* \;\approx\; \frac{\ln n}{\ln k} \qquad\text{(turnaj velikosti } k)
\]
\pop{$\tau^*$ = počet generací, za který jediná kopie nejlepšího jedince (bez variace) zaplní
celou populaci --- kvantifikace \emph{selekčního tlaku}. Pro fitness-proporcionální selekci je
doba převzetí delší a~závisí na rozložení fitness.}

\fml{Generation gap a~modely populace}
\[
G = \frac{\lambda}{n}
\]
\pop{Podíl nahrazované populace: generační model $G=1$ (celá populace nahrazena potomky),
steady-state $\lambda \in \{1,2\}$. \emph{Elitismus}: $k$ nejlepších se bezpodmínečně kopíruje
dál (zaručuje monotonii nejlepší fitness, je nutnou podmínkou řady konvergenčních vět).
Deterministická selekce ES: $(\mu+\lambda)$ vybírá $\mu$ nejlepších z~rodičů \emph{i}~potomků,
$(\mu,\lambda)$ jen z~potomků.}

\podkap{2}{Genetické algoritmy}

\podpodkap{2}{Binární reprezentace}

\fml{Binární kódování reálného čísla}
\[
x \;=\; a \;+\; z \cdot \frac{b-a}{2^k - 1},
\qquad \text{přesnost} = \frac{b-a}{2^k-1}
\]
\pop{$x \in [a,b]$ kódované na $k$ bitů jako celé číslo $z \in \{0,\dots,2^k-1\}$.
Výhoda: explicitní kontrola nad přesností. Nevýhoda: \emph{Hammingovy útesy} (sousední čísla
se liší v~mnoha bitech, např.\ $0111 \to 1000$) --- řeší se \emph{Grayovým kódem}.}

\podpodkap{4}{Mutace a~inverze}

\fml{Bitová mutace}
\[
p_M \;\approx\; \frac{1}{m}
\]
\pop{$m$ je délka řetězce; každý bit se převrací nezávisle s~pravděpodobností $p_M$; doporučená hodnota znamená
\uv{v~průměru se změní jeden bit na jedince}. Typické další parametry SGA:
$n \in [50,500]$, $p_C \in [0{,}6;\,0{,}9]$.}

\podpodkap{5}{Škálování fitness}

\fml{Lineární škálování}
\[
f' = a f + b, \qquad \bar{f'} = \bar{f}, \qquad
f'_{\max} = C_{\mathrm{mult}} \cdot \bar{f}, \quad C_{\mathrm{mult}} \in [1{,}2;\,2]
\]
\pop{Dvě podmínky určují $a$ a~$b$. Zachovává průměr a~omezuje náskok nejlepšího jedince na
$C_{\mathrm{mult}}$-násobek průměru. Nutná pojistka proti záporným hodnotám u~nejhorších.}

\fml{Sigma truncation, mocninné a~Boltzmannovo škálování}
\[
f' = \max\big(0,\; f - (\bar f - c\,\sigma_f)\big), \qquad
f' = f^{\,k}, \qquad
f' = e^{f/T}
\]
\pop{$\sigma_f$ je směrodatná odchylka fitness, $c \in [1;3]$, $T$ teplota. Sigma truncation odstraní
odlehlé nízké hodnoty a~automaticky se přizpůsobuje rozptylu; u~mocninného škálování $k$
v~průběhu běhu roste (tlak roste). Nejrobustnější a~dnes nejobvyklejší je ale
\emph{pořadové} škálování (nahrazení fitness pořadím).}

\podkap{3}{Reprezentace jedince a~genetické operátory}

\podpodkap{2}{Reálná reprezentace}

\fml{Aritmetické křížení}
\[
\vec z \;=\; \alpha \vec x + (1-\alpha)\vec y, \qquad \alpha \in (0,1)
\]
\pop{$\vec x, \vec y$ jsou rodiče a~$\vec z$ potomek --- konvexní kombinace rodičů; pro $\alpha = 1/2$ prostý průměr. Varianty:
\emph{jednoduché} (všechny složky), \emph{jednoduché aritmetické} (jedna náhodná složka),
\emph{celé/single} (kombinace s~jednobodovým křížením). Nevýhoda: potomci leží vždy
\emph{uvnitř} konvexního obalu rodičů $\Rightarrow$ populace se stahuje.}

\fml{BLX-$\alpha$ (\emph{blend crossover})}
\[
z_i \sim U\big[\min(x_i,y_i) - \alpha I,\ \max(x_i,y_i) + \alpha I\big],
\qquad I = |x_i - y_i|,\ \ \alpha \approx 0{,}5
\]
\pop{Rozšiřuje interval o~$\alpha I$ na obě strany, takže potomci mohou ležet \emph{i~vně}
rodičů --- odstraňuje stahování populace, které trápí čistě aritmetické křížení.}

\fml{SBX (\emph{simulated binary crossover})}
\[
z_{1,i} = \tfrac12\big[(1+\beta)x_i + (1-\beta)y_i\big],\qquad
z_{2,i} = \tfrac12\big[(1-\beta)x_i + (1+\beta)y_i\big]
\]
\[
\beta = (2u)^{1/(\eta+1)} \ \text{ pro } u \le \tfrac12, \qquad
\beta = \big(\tfrac{1}{2(1-u)}\big)^{1/(\eta+1)} \ \text{ pro } u > \tfrac12,
\qquad u \sim U(0,1)
\]
\pop{Napodobuje chování jednobodového binárního křížení v~reálném prostoru: potomci vzniknou
\uv{roztažením} dvojice rodičů kolem jejich \emph{středu} faktorem $\beta$. Platí
$|z_1 - z_2| = \beta\,|x-y|$ a~střed rodičů zůstává zachován. Parametr $\eta$ (distribuční
index) řídí, jak blízko rodičům potomci padnou --- velké $\eta$ $\Rightarrow$ blízko; $u \sim U(0,1)$ se losuje pro každou složku.
Standardní operátor v~NSGA-II.}

\fml{Gaussovská mutace reálného vektoru}
\[
x_i' = x_i + \sigma \cdot N_i(0,1)
\]
\pop{$\sigma$ je velikost kroku (\emph{step size}); její řízení je hlavní téma evolučních
strategií.}

\podkap{4}{Genetické a~evoluční programování}

\podpodkap{3}{Bloat}

\fml{Parsimony pressure (boj s~bloatem v~GP)}
\[
f' = f - \alpha \cdot \text{velikost}
\]
\pop{\emph{Bloat} = tendence programů růst bez odpovídajícího zlepšení fitness (hromadění
\emph{intronů} typu \texttt{(+ x 0)}). Alternativy: tvrdé limity na hloubku/počet uzlů,
lexikografické porovnání (při shodné fitness vyhrává menší jedinec), anti-bloat operátory
(\emph{hoist}, \emph{shrink}), nebo velikost jako druhé kritérium ve~vícekriteriální
optimalizaci.}

\kap{2}{Teorie schémat a~pravděpodobnostní modely SGA}

\podkap{1}{Teorie schémat}

\podpodkap{1}{Schéma, jeho řád a~definující délka}

\fml{Schéma, řád, definující délka}
\[
S \in \{0,1,*\}^m, \qquad o(S) = \text{počet pevných pozic}, \qquad
d(S) = \text{poslední pevná} - \text{první pevná pozice}
\]
\pop{Schéma s~$r$ hvězdičkami reprezentuje $2^r$ jedinců. Příklad: $S = 1{*}{*}0{*}1$ má
$m=6$, $o(S)=3$, $d(S) = 6-1 = 5$. $F(S,t)$ = průměrná fitness reprezentantů $S$
v~populaci $P(t)$ --- pozor, závisí na \emph{aktuální} populaci, nejen na $S$ a~$f$.}

\fml{Kombinatorika schémat}
\[
\#\{\text{schémat délky } m\} = 3^m, \qquad
\#\{\text{schémat, jejichž reprezentantem je jeden jedinec}\} = 2^m
\]
\pop{Populace $n$ jedinců reprezentuje mezi $2^m$ a~$n\cdot 2^m$ schématy (podle míry jejich
podobnosti). Odtud \emph{implicitní paralelismus}: GA pracuje explicitně s~$n$ jedinci, ale
implicitně zpracovává mnohem víc schémat --- užitečně zpracovávaných je podle Hollanda
řádu $n^3$.}

\podpodkap{2}{Hollandova věta o~schématech}

\fml{Věta o~schématech (Holland, 1975)}
\[
\mathbb{E}\big[C(S,t+1)\big] \;\ge\; C(S,t)\,\frac{F(S,t)}{\bar f(t)}
\left[\,1 - p_C \frac{d(S)}{m-1} - p_M\, o(S)\right]
\]
\pop{$C(S,t)$ = počet reprezentantů $S$ v~$P(t)$. Slovně: \textbf{krátká} (malé $d(S)$),
\textbf{nízkého řádu} (malé $o(S)$) a~\textbf{nadprůměrná} ($F(S,t) > \bar f(t)$) schémata se
v~po sobě jdoucích generacích množí exponenciálně. To je \emph{hypotéza stavebních bloků}:
GA skládá řešení z~krátkých, dobrých stavebních bloků.}

\fml{Odvození, krok 1 --- selekce}
\[
p_s(S) = \sum_{v \in S \cap P(t)} \frac{f(v)}{F(t)} = \frac{C(S,t)\,F(S,t)}{F(t)},
\qquad
\mathbb{E}\big[C(S,t+1)\big] = n\,p_s(S) = C(S,t)\,\frac{F(S,t)}{\bar f(t)}
\]
\pop{$F(t) = \sum_{u \in P(t)} f(u)$, $\bar f(t) = F(t)/n$. Je-li schéma trvale nadprůměrné
o~$\varepsilon$, tj.\ $F(S,t) = \bar f(t)(1+\varepsilon)$, dostáváme
$C(S,t) = C(S,0)\,(1+\varepsilon)^t$ --- \textbf{exponenciální růst} (pokles pro
$\varepsilon<0$).}

\fml{Odvození, kroky 2 a~3 --- křížení a~mutace}
\[
p_{\mathrm{surv}}^{\mathrm{cross}}(S) \;\ge\; 1 - p_C\,\frac{d(S)}{m-1},
\qquad
p_{\mathrm{surv}}^{\mathrm{mut}}(S) = (1-p_M)^{o(S)} \;\approx\; 1 - p_M\,o(S)
\]
\pop{Jednobodové křížení volí bod řezu z~$m-1$ možností; schéma je jistě rozbito, padne-li řez
mezi první a~poslední pevnou pozici. Nerovnost \uv{$\ge$} platí proto, že i~při řezu uvnitř
může být schéma obnoveno druhým rodičem. Mutace: schéma přežije, nezmění-li se \emph{žádná}
z~$o(S)$ pevných pozic; aproximace platí pro $p_M \ll 1$. Složením obou (a~zanedbáním
součinu) vzniká věta o~schématech.}

\podpodkap{5}{Kritika teorie schémat}

\fml{Deceptivní úloha (kdy GA selhává)}
\[
S_1 = (111{*}\ldots{*}), \quad S_2 = ({*}\ldots{*}11)
\quad\text{nadprůměrná, ale}\quad
f(111{*}{*}{*}{*}{*}11) \;\ll\; f(000{*}{*}{*}{*}{*}00)
\]
\pop{Selekce vede GA k~jedničkovým blokům, přestože optimum leží jinde --- stavební bloky
\uv{lžou}. Hlavní protipříklad k~hypotéze stavebních bloků.}

\podkap{2}{Pravděpodobnostní modely jednoduchého GA}

\podpodkap{3}{Operátor $\mathcal{G}$}

\fml{Operátor $\mathcal{G}$ a~jeho rozklad}
\[
p(t+1) = \mathcal{G}\big(p(t)\big), \qquad \mathcal{G} = \mathcal{M} \circ \mathcal{F}
\]
\pop{$p \in \Lambda$ (simplex) je vektor podílů jednotlivých řetězců v~populaci,
$s$ vektor selekčních pravděpodobností. $\mathcal{F}$ odpovídá selekci, $\mathcal{M}$
\emph{míchání} (křížení $+$ mutace dohromady). Iterováním $\mathcal{G}$ z~počáteční populace
známe úplné chování SGA nad \emph{nekonečnou} populací.}

\fml{Část $\mathcal{F}$ --- fitness-proporcionální selekce maticově}
\[
s(t) \;=\; \frac{\mathbf{F}\,p(t)}{\sum_j \mathbf{F}_{jj}\,p_j(t)}
\]
\pop{$\mathbf F$ je diagonální matice s~$\mathbf F_{ii} = f(i)$. Čitatel $(\mathbf Fp)_i =
f(i)p_i$ je \uv{váha} řetězce $i$, jmenovatel je průměrná fitness populace. Podíl řetězce se
tedy mění přesně o~faktor $f(i)/\bar f$ --- v~souladu s~větou o~schématech.}

\fml{Část $\mathcal{M}$ --- míchání}
\[
r(i,j,k) = \Pr[\text{potomek } k \text{ vznikne z~rodičů } i,j],
\qquad
\mathbb{E}\big[p_k(t+1)\big] = \sum_{i}\sum_{j} s_i(t)\, s_j(t)\, r(i,j,k)
\]
\pop{Jde o~\emph{kvadratický} operátor na simplexu. Klíčové zjednodušení: jednobodové křížení
i~bitová mutace \uv{nevidí} absolutní hodnoty bitů, jen shody a~neshody, takže
$r(i,j,k) = r(i \oplus k,\ j \oplus k,\ 0)$ ($\oplus$ = XOR). Stačí tedy jediná
\emph{míchací matice} $\mathbf M_{ij} = r(i,j,0)$ rozměru $2^l \times 2^l$.}

\podpodkap{5}{Konečné populace: Markovův řetězec}

\fml{Konečné populace jako Markovův řetězec}
\[
N \;=\; \binom{n + 2^l - 1}{\,2^l - 1\,},
\qquad
\mathbf{Q}_{i,j} \;=\; n! \prod_{y=0}^{2^l-1}
\frac{\big(\mathcal{G}(p_i)_y\big)^{\mathbf{Z}(y,j)}}{\mathbf{Z}(y,j)!}
\]
\pop{Stavy řetězce = jednotlivé možné populace (multimnožiny $n$ řetězců délky $l$), jejich
počet je $N$. $\mathbf Z(y,i)$ = počet výskytů jedince $y$ v~populaci $i$. Přechod je
\emph{multinomické} losování $n$ jedinců podle rozdělení $\mathcal G(p)$. Model je exaktní,
ale $N$ roste tak rychle, že je prakticky nepoužitelný.}

\kap{3}{Evoluční strategie, diferenciální evoluce, koevoluce, otevřená evoluce}

\podkap{1}{Evoluční strategie}

\podpodkap{2}{$(1+1)$-ES a~pravidlo 1/5}

\fml{Pravidlo 1/5 úspěchu (Rechenberg)}
\[
p > 1/5 \Rightarrow \sigma \gets \sigma / c, \qquad
p < 1/5 \Rightarrow \sigma \gets \sigma \cdot c, \qquad
0{,}8 \le c \le 1
\]
\pop{$p$ = podíl úspěšných mutací za posledních $k$ iterací, $c$ = multiplikativní konstanta. \emph{Optimální poměr úspěšných
mutací je přibližně $1/5$}: je-li úspěchů víc, je krok příliš malý (zvětši $\sigma$),
je-li jich méně, je krok příliš velký (zmenši $\sigma$). Jde o~\emph{adaptivní}, nikoli
samoadaptivní řízení parametru --- $\sigma$ není součástí genomu.}

\podpodkap{3}{Samoadaptace strategických parametrů}

\fml{Samoadaptace strategických parametrů (Schwefelova pravidla)}
\[
\sigma_i' = \sigma_i \cdot \exp\big(\tau' N(0,1) + \tau N_i(0,1)\big),
\qquad
x_i' = x_i + \sigma_i' \cdot N_i(0,1)
\]
\[
\tau' \propto \frac{1}{\sqrt{2n}}, \qquad \tau \propto \frac{1}{\sqrt{2\sqrt{n}}}
\]
\pop{Jedinec má tvar $C = [\vec x, \vec\sigma]$ --- strategické parametry jsou
\textbf{součástí genomu} a~podléhají mutaci i~selekci. $N(0,1)$ je \emph{jedno} losování
společné pro celý jedinec, $N_i(0,1)$ se losuje pro každou složku zvlášť. Exponenciála
zaručuje kladnost $\sigma$ a~symetrii zvětšení/zmenšení; $\tau'$ a~$\tau$ jsou rychlosti učení (globální a~pro jednotlivé složky).
\textbf{Pořadí je zásadní: nejprve se mutuje $\sigma$, teprve pak se s~novým $\sigma$
mutuje $x$} --- jinak by se hodnotilo staré $\sigma$.}

\fml{Kovarianční matice z~$\sigma_i$ a~úhlů natočení}
\[
c_{ii} = \sigma_i^2, \qquad
c_{ij} = \tfrac12(\sigma_i^2 - \sigma_j^2)\tan(2\alpha_{ij})
\]
\pop{Tři úrovně samoadaptace: jedno $\sigma$ (izotropní mutace), $n$ hodnot $\sigma_i$
(osově orientovaný elipsoid), $\sigma_i$ + úhly $\alpha_{ij}$ (plná kovarianční matice,
libovolně natočený elipsoid). CMA-ES místo losování $\alpha_{ij}$ odhaduje $\mathbf C$
z~\emph{cesty evoluce} úspěšných kroků.}

\podkap{2}{Diferenciální evoluce}

\podpodkap{2}{Algoritmus DE/rand/1/bin}

\fml{Mutace --- dárcovský vektor}
\[
\vec v_i \;=\; \vec x_{a} \;+\; F\,(\vec x_{b} - \vec x_{c}),
\qquad F \in [0,2],\ \text{typicky } F \approx 0{,}8
\]
\pop{$\vec x_a, \vec x_b, \vec x_c$ jsou tři navzájem různí jedinci, různí od $\vec x_i$; $F$ je mutační konstanta (škálovací faktor).
Klíčová myšlenka: \emph{velikost kroku se bere z~populace samotné} --- na začátku (rozptýlená
populace) jsou rozdíly velké, ke konci malé. Adaptace měřítka je tedy zdarma, bez
strategických parametrů.}

\fml{Binomické křížení --- zkušební vektor}
\[
u_{j,i} =
\begin{cases}
v_{j,i} & \text{pokud } \mathrm{rand}_j \le C \text{ nebo } j = I_{\mathrm{rand}},\\[2pt]
x_{j,i} & \text{jinak}
\end{cases}
\]
\pop{$\mathrm{rand}_j \sim U(0,1)$, $C \in [0,1]$ je práh křížení, $I_{\mathrm{rand}}$
náhodný index z~$\{1,\dots,D\}$. Podmínka $j = I_{\mathrm{rand}}$ je \uv{pojistka}:
zaručuje, že alespoň jedna složka pochází od dárce, takže potomek není nikdy identickou kopií
rodiče.}

\fml{Environmentální selekce (souboj 1:1)}
\[
\vec x_i(t+1) =
\begin{cases}
\vec u_i & \text{pokud } f(\vec u_i) \le f(\vec x_i(t)),\\
\vec x_i(t) & \text{jinak}
\end{cases}
\]
\pop{Selekce je elitistická na úrovni jednotlivce --- populace se nikdy nezhorší.
Rodičovská selekce je triviální: rodičem je postupně \emph{každý} jedinec.
Varianty mutace se značí \texttt{DE/}\emph{báze}\texttt{/}\emph{počet diferencí}\texttt{/}\emph{křížení},
např.\ \texttt{DE/best/1/bin} používá místo $\vec x_a$ nejlepšího jedince.}

\podkap{3}{Koevoluce}

\podpodkap{5}{Evoluce kooperace a~vězňovo dilema}

\fml{Vězňovo dilema jako model evoluce kooperace}
\[
\begin{array}{c|cc}
 & C & D\\\hline
C & R/R & S/T\\
D & T/S & P/P
\end{array}
\qquad
T > R > P > S, \quad 2R > T+S
\]
\pop{Typicky $R=-1$, $T=0$, $P=-2$, $S=-3$. Zrada je dominantní strategie, $DD$ je Nashovo
ekvilibrium, ale $CC$ je Pareto-lepší. V~\emph{iterované} verzi (Axelrodovy turnaje) vítězí
\emph{tit for tat}: hodná (nezradí první), oplácející, odpouštějící a~průhledná strategie.
Formální model recipročního altruismu --- kooperace se vyplatí při dostatečném
\uv{stínu budoucnosti}.}

\podpodkap{6}{Diverzita, druhy a~niky}

\fml{Fitness sharing}
\[
f'(i) = \frac{f(i)}{\sum_j sh\big(d(i,j)\big)},
\qquad
sh(d) = \begin{cases}
1-\big(d/\sigma_{\mathrm{share}}\big)^{\alpha} & d < \sigma_{\mathrm{share}},\\
0 & \text{jinak}
\end{cases}
\]
\pop{Fitness se \uv{dělí} mezi podobné jedince, takže obsazená nika se stává méně
atraktivní --- vzniká tlak na pokrytí \emph{různých} nik. $\sigma_{\mathrm{share}}$ je
poloměr niky, $d$ vzdálenost (Hammingova, fenotypová, \dots), $\alpha$ exponent sdílecí funkce $sh$. Problém je vždy volba
$\sigma_{\mathrm{share}}$; NSGA-II ji proto nahradilo \emph{crowding distance}.}

\podkap{4}{Otevřená evoluce}

\fml{Novelty search --- řídkost okolí}
\[
\rho(x) \;=\; \frac{1}{k}\sum_{i=1}^{k} \mathrm{dist}\big(x, \mu_i\big)
\]
\pop{$\mu_i$ jsou $k$ nejbližších sousedů $x$ v~aktuální populaci \textbf{a~v~permanentním
archivu}, měřeno v~\emph{prostoru chování} (např.\ koncová pozice robota v~bludišti).
Fitness $= \rho$; původní účelová funkce se \emph{nepoužívá vůbec}. Motivace (Lehman \&
Stanley): u~deceptivních úloh je cílová funkce zavádějící a~vede do lokálních optim.}

\kap{4}{Rojové optimalizační algoritmy}

\podkap{1}{Optimalizace rojem částic (PSO)}

\fml{PSO --- pohybové rovnice}
\[
\vec v_i \;\gets\; \underbrace{w\,\vec v_i}_{\text{setrvačnost}}
\;+\; \underbrace{c_1 r_1 (\vec p_i - \vec x_i)}_{\text{kognitivní složka}}
\;+\; \underbrace{c_2 r_2 (\vec g - \vec x_i)}_{\text{sociální složka}},
\qquad
\vec x_i \;\gets\; \vec x_i + \vec v_i
\]
\pop{$\vec p_i$ (\emph{pBest}) = nejlepší pozice, kterou částice $i$ kdy navštívila;
$\vec g$ (\emph{gBest}) = nejlepší pozice nalezená rojem. $r_1, r_2 \sim U(0,1)$ se losují
pro každou složku zvlášť, $c_1, c_2$ jsou koeficienty učení (původně $c_1=c_2=2$), $w$ je
setrvačnost. Částice pravidelně \emph{přeletí} přes $\vec g$ --- to je zdroj explorace;
proto se rychlost omezuje ($|v_j| \le v_{\max}$) nebo se používá konstrikční faktor.}

\podkap{2}{Optimalizace mravenčí kolonií (ACO)}

\fml{ACO --- pravidlo přechodu (Ant System)}
\[
w_{ij} \;=\; \big[\tau_{ij}\big]^{\alpha}\,\big[\eta_{ij}\big]^{\beta},
\qquad
p_{ij}^{k} \;=\; \frac{w_{ij}}{\displaystyle\sum_{l \in \mathcal{N}_i^k} w_{il}},
\qquad \eta_{ij} = \frac{1}{d_{ij}}
\qquad\Longrightarrow\qquad
p_j \;\propto\; \frac{\tau_j^{\,\alpha}}{d_j^{\,\beta}}
\]
\pop{$\tau_{ij}$ = feromon na hraně, $\eta_{ij}$ = heuristická \uv{viditelnost} (převrácená
vzdálenost), $\mathcal{N}_i^k$ = dosud nenavštívená města mravence $k$, $w_{ij}$ = skóre
kandidáta. Je to \textbf{obyčejná ruleta}: skóre / součet skóre. Indexy $i,k$ jsou během
jednoho rozhodnutí konstanty, takže si stačí pamatovat pravý tvar --- \uv{feromon nahoře,
vzdálenost dole} --- a~dopsat normalizaci.
$\alpha,\beta$ váží feromon vs.\ heuristiku (typicky $\alpha=1$, $\beta \in [2,5]$);
$\alpha=0$ dá hladový algoritmus, $\beta=0$ čistě feromonové hledání (stagnuje).}

\fml{ACO --- aktualizace feromonu}
\[
\tau_{ij} \;\gets\; (1-\rho)\,\tau_{ij} \;+\; \sum_{k=1}^{m} \Delta\tau_{ij}^{k},
\qquad
\Delta\tau_{ij}^{k} =
\begin{cases}
Q/L_k & \text{prošel-li mravenec } k \text{ hranou } (i,j),\\
0 & \text{jinak}
\end{cases}
\]
\pop{$\rho \in (0,1)$ je míra odpařování (zapomínání --- brání předčasné konvergenci),
$L_k$ délka trasy mravence $k$, $Q$ konstanta, $m$ počet mravenců, $\Delta\tau^k_{ij}$ příspěvek mravence $k$. \textbf{Kratší trasa $\Rightarrow$ více
feromonu.} ACS navíc odpařuje lokálně během konstrukce a~ukládá jen po nejlepší trase.}

\kap{5}{Lokální prohledávání a~memetické algoritmy}

\podkap{2}{Simulované žíhání}

\fml{Metropolisovo kritérium (simulované žíhání)}
\[
P(\text{přijmout } y) =
\begin{cases}
1, & \Delta \le 0 \quad (\text{zlepšení}),\\[2pt]
e^{-\Delta / T}, & \Delta > 0 \quad (\text{zhoršení}),
\end{cases}
\qquad \Delta = f(y) - f(x)
\]
\pop{Pro minimalizaci, kandidát $y \in N(x)$. $T$ je teplota: při vysokém $T$ se algoritmus
chová jako náhodná procházka (např.\ $\Delta=2$, $T=5$ dá $P=0{,}67$), při nízkém jako čistý
horolezec ($T=0{,}5$ dá $P=0{,}018$). Při pevné $T$ konverguje k~Boltzmannovu rozdělení
$\pi(x) \propto e^{-f(x)/T}$.}

\fml{Rozvrhy chlazení}
\[
T_{k+1} = \alpha T_k,\ \ \alpha \in [0{,}8;\,0{,}99]
\qquad
T_k = T_0 - \beta k
\qquad
T_k = \frac{c}{\log(k+1)}
\]
\pop{$T_k$ je teplota v~$k$-té iteraci, $\alpha$ chladicí faktor, $\beta$ krok lineárního poklesu. Geometrický (v~praxi nejběžnější), lineární, logaritmický. \textbf{Jen pro logaritmický
je dokázána konvergence} k~množině globálních optim s~pravděpodobností 1 (při dostatečně
velkém $c$, souvisejícím s~výškou nejhlubší bariéry) --- prakticky je ale nepoužitelně
pomalý.}

\kap{6}{Aplikace evolučních algoritmů}

\podkap{1}{Evoluce pravidlových a~expertních systémů}

\podpodkap{2}{Learning classifier systems (Michigan)}

\fml{XCS --- přesnost klasifikátoru}
\[
\kappa = \alpha\,(\epsilon/\epsilon_0)^{-\nu} \quad (\epsilon > \epsilon_0), \qquad
\kappa = 1 \quad\text{jinak}
\]
\pop{Klasifikátor je pětice $(c, a, p, \epsilon, F)$: podmínka, akce, predikce odměny,
chyba predikce, fitness. Wilsonova myšlenka: \emph{predikce} má odhadovat odměnu, ale
\emph{evoluce} má preferovat \textbf{spolehlivé (přesné)} klasifikátory --- fitness je proto
relativní přesnost v~rámci action setu, ne síla. $\kappa$ je přesnost, $\alpha, \epsilon_0, \nu$ parametry (práh a~strmost poklesu). $p$ a~$\epsilon$ se aktualizují pravidly
Q-learningu ($r + \gamma\max_a p_a$), GA běží \emph{nikově} jen uvnitř action setu.}

\fml{ZCS --- posílení}
\[
\text{každému v } A:\ +\frac{\beta r}{|A|},
\qquad
\text{každému v } A_{-1}:\ +\frac{\gamma B}{|A_{-1}|}
\]
\pop{$A$ = action set (pravidla prosazující zvolenou akci), $A_{-1}$ = předchozí action set,
$B$ = \uv{kbelík} (podíl $\beta$ odebraný ze síly pravidel v~$A$), $r$ = odměna z~prostředí.
Zjednodušení Hollandovy \emph{bucket brigade}, kde pravidlo muselo \uv{zaplatit} tomu, které
ho aktivovalo.}

\podpodkap{4}{Připomenutí: metriky klasifikace}

\fml{Metriky klasifikace (fitness pravidlových systémů)}
\[
\text{přesnost} = \frac{TP+TN}{TP+TN+FP+FN}, \quad
\text{senzitivita} = \frac{TP}{TP+FN}, \quad
\text{specificita} = \frac{TN}{TN+FP}
\]
\pop{Fitness pravidlového systému (LCS, GIL) bývá kombinací přesnosti, pokrytí
a~jednoduchosti (počet pravidel) --- tedy přirozeně vícekriteriální úloha.
Podrobněji viz odd.~\ref{ssec:matice-zamen}.}

\podkap{2}{Neuroevoluce}

\podpodkap{3}{NEAT}

\fml{NEAT --- míra podobnosti genomů (speciace)}
\[
\delta \;=\; \frac{c_1 E}{N} + \frac{c_2 D}{N} + c_3 \overline{\Delta W}
\]
\pop{\textbf{Předpoklad:} síť je v~NEAT \emph{seznam hran} a~každá hrana nese
\emph{inovační číslo} --- pořadové číslo přidělené při prvním vzniku dané strukturální
novinky. Hrany se stejným číslem tedy mají \textbf{společný evoluční původ}. Porovnání dvou
genomů jen roztřídí jejich hrany do tří skupin: \emph{shodné} (totéž inovační číslo v~obou),
\textbf{disjoint} (chybí u~druhého, ale \emph{uvnitř} rozsahu jeho čísel) a~\textbf{excess}
(leží \emph{za} nejvyšším inovačním číslem druhého genomu).

\emph{Příklad:} $A=\{1,2,3,4,5,8\}$, $B=\{1,2,3,6,7\}$. Shodné jsou 1,\,2,\,3.
Nejvyšší číslo $B$ je 7, takže geny 4,\,5 (u~$A$) a~6,\,7 (u~$B$) leží uvnitř rozsahu
$\Rightarrow D=4$; gen 8 je za ním $\Rightarrow E=1$; $N=6$ (větší genom).

\textbf{Co který člen dělá:} $E$ a~$D$ měří rozdíl \emph{struktury} (kolik hran mají genomy
navíc jeden proti druhému), $\overline{\Delta W}$ rozdíl \emph{parametrů} (průměrný rozdíl vah
přes \emph{shodné} geny --- jen tam má porovnávat vah smysl). $N$ = velikost většího genomu
normalizuje, aby velké sítě nevycházely vzdálené jen svou velikostí; $c_1,c_2,c_3$ jsou váhy,
kterými volíš, zda víc rozhoduje topologie, nebo váhy.
Sítě s~$\delta$ pod prahem tvoří jeden \textbf{druh} a~sdílejí fitness --- nová topologie tak
dostane čas doladit váhy dřív, než ji vytlačí zavedená řešení.}

\podkap{3}{Kombinatorické úlohy a~práce s~omezeními}

\podpodkap{2}{Tři obecné strategie pro omezení}

\fml{Penalizace omezení (problém batohu)}
\[
f'(x) = \sum_i x_i v_i - \alpha \max\Big(0, \sum_i x_i c_i - C_{\max}\Big)
\]
\pop{$C_{\max}$ je kapacita batohu, $v_i$ cena a~$c_i$ objem předmětu $i$, $x_i \in \{0,1\}$ jeho výběr. Tři obecné strategie pro omezení: \emph{(1)}~\textbf{penalizace} (uvedený vzorec ---
volba $\alpha$ je zásadní: malá $\Rightarrow$ populace se plní nepřípustnými řešeními, velká
$\Rightarrow$ nelze projít nepřípustnou oblastí ke slibnému řešení); \emph{(2)}~\textbf{oprava}
(u~batohu vyhazuj předměty s~nejhorším poměrem $v_i/c_i$ --- baldwinovsky jen pro fitness,
lamarckovsky i~do genotypu); \emph{(3)}~\textbf{dekodér}, kde každý genotyp reprezentuje
přípustné řešení. Čtvrtá možnost: porušení omezení jako druhé kritérium.}

\podkap{4}{Vícekriteriální optimalizace}

\podpodkap{1}{Dominance a~Pareto fronta}

\fml{Dominance a~Pareto optimalita}
\[
x \preceq y \iff \forall i:\ f_i(x) \le f_i(y),
\qquad
x \prec y \iff x \preceq y \ \wedge\ \exists j:\ f_j(x) < f_j(y)
\]
\pop{(Vše minimalizujeme.) $\prec$ je \emph{ostré částečné} uspořádání --- ireflexivní
a~tranzitivní, ale \textbf{nikoli úplné}, takže existuje celá množina navzájem
neporovnatelných optim. Řešení nedominované žádným přípustným je \emph{Pareto-optimální};
množina takových řešení je \emph{Pareto množina}, její obraz v~prostoru kritérií
\textbf{Pareto fronta}.}

\podpodkap{2}{Skalarizace --- jednoduchá řešení}

\fml{Skalarizace}
\[
f = \sum_i w_i f_i
\qquad
\min f_1 \text{ s.t. } f_i \le \varepsilon_i
\qquad
\min \max_i \lambda_i \,|f_i(x) - z_i^*|
\]
\pop{$w_i$ a~$\lambda_i$ jsou váhy kritérií. Lineární skalarizace, $\varepsilon$-constraint, Čebyševova skalarizace ($z^*$ =
referenční/ideální bod). \textbf{Klíčová nevýhoda lineární skalarizace: nekonvexní části
fronty nelze nikdy nalézt}; $\varepsilon$-constraint i~Čebyševova je najdou. Každý běh navíc
dá jen \emph{jeden} bod fronty.}

\podpodkap{4}{NSGA-II}

\fml{NSGA-II --- crowding distance}
\[
d_i \;=\; \sum_{m=1}^{k}
\frac{f_m(i+1) - f_m(i-1)}{f_m^{\max} - f_m^{\min}}
\]
\pop{Počítá se \emph{pro každou frontu zvlášť}: pro každé kritérium $m$ se jedinci seřadí
podle $f_m$, krajní dostanou $\infty$ a~ostatním se přičte normalizovaná vzdálenost sousedů.
Vyjadřuje \uv{obvod kvádru} tvořeného nejbližšími sousedy --- \emph{velká} hodnota znamená
řídce obsazenou oblast, kterou chceme zachovat. Nahrazuje $\sigma_{\mathrm{share}}$
původního NSGA (žádný parametr navíc). Selekce: nejprve nižší fronta, při shodě větší $d_i$.}

\podpodkap{6}{Metriky kvality aproximace}

\fml{Hypervolume ($S$-metrika)}
\[
HV = \mathrm{vol}\Big(\bigcup_{x \in \text{fronta}} [f(x), r]\Big)
\]
\pop{Objem oblasti dominované nalezenou frontou a~ohraničené \emph{referenčním bodem} $r$
(zvoleným tak, aby byl dominován všemi řešeními). Jediná běžná metrika \emph{striktně
monotónní vůči dominanci}, která navíc \textbf{nevyžaduje znalost skutečné Pareto fronty}.
Příklad: fronta $\{(2,7),(4,5),(8,2)\}$, $r=(10,10)$; po svislých pruzích
$2\cdot8 + 4\cdot5 + 2\cdot3 = 16+20+6 = 42$. GD/IGD a~spread naopak skutečnou frontu
znát potřebují.}

%% ========================================================================
\okruh{03}{Dobývání znalostí}{Příprava dat a~relevance atributů, asociační pravidla, klasifikace, shlukování, vyhodnocování modelů, analýza sociálních sítí.}
%% ========================================================================

\kap{2}{Příprava dat, výběr atributů a~analýza jejich relevance}

\podkap{3}{Čištění dat: šum a~odlehlé hodnoty}

\fml{Odlehlá hodnota podle mezikvartilového rozpětí}
\[
\mathrm{IQR}=Q_3-Q_1,
\qquad
x > Q_3 + 1{,}5\cdot \mathrm{IQR} \quad\text{nebo}\quad x < Q_1 - 1{,}5\cdot \mathrm{IQR}
\]
\pop{$Q_1,Q_2,Q_3$ jsou kvartily (25~\%, 50~\% = medián, 75~\% pozorování pod nimi).
Standardní pravidlo pro detekci \emph{outlieru}; základ krabicového grafu.}

\podkap{5}{Transformace a~standardizace atributů}

\fml{Min-max normalizace (\emph{range standardization})}
\[
\mathrm{range}(x_{if})=\frac{x_{if}-\min(f)}{\max(f)-\min(f)} \in [0,1],
\qquad
\mathrm{range}_{[-1,1]} = 2\cdot\mathrm{range}(x_{if})-1
\]
\pop{$f$ je atribut, $x_{if}$ jeho hodnota u~$i$-tého objektu. \textbf{Proč:} bez
standardizace ovládne vzdálenost atribut s~největším rozsahem --- pro
$x_i=(0{,}1;\,20)$ a~$x_j=(0{,}9;\,720)$ je $\mathrm{dist} = 700{,}0005$, tedy první atribut
nemá prakticky žádný vliv. \emph{Desítkové škálování} dělí nejmenší mocninou 10, která
udrží hodnoty v~$[-1,1]$.}

\fml{Standardizace podle průměrné absolutní odchylky}
\[
m_f=\tfrac1n\big(x_{1f}+\cdots+x_{nf}\big),\qquad
s_f=\tfrac1n\big(|x_{1f}-m_f|+\cdots+|x_{nf}-m_f|\big),\qquad
z(x_{if})=\frac{x_{if}-m_f}{s_f}
\]
\pop{Dává nulový průměr a~jednotkovou \emph{střední absolutní} odchylku. Oproti klasickému
$z$-skóre (se směrodatnou odchylkou) je robustnější vůči odlehlým hodnotám, protože se
odchylky neumocňují.}

\podkap{6}{Diskretizace numerických atributů}

\podpodkap{1}{Fuzzy diskretizace}

\fml{Fuzzy diskretizace --- charakteristické funkce a~váha objektu}
\[
\mu_{i-1}(x)=\frac{x_{i+1}-x}{x_{i+1}-x_{i-1}},\qquad
\mu_{i+1}(x)=\frac{x-x_{i-1}}{x_{i+1}-x_{i-1}},\qquad
w = \prod_j \mu_j < 1
\]
\pop{Na přechodové zóně mezi dvěma intervaly charakteristické funkce lineárně klesají, resp.\
rostou, a~jejich součet je 1. Jedné numerické hodnotě tak mohou odpovídat \emph{dvě} sousední
kategorie; objekt v~zóně vstupuje do dat jako \uv{fuzzy objekt} s~vahou $w$ danou součinem
přes všechny atributy. Motivace: ostrá diskretizace ztrácí informaci a~zavádí
\emph{kontradikce}.}

\podkap{8}{Analýza relevance atributů: filtry a~wrappery}

\podpodkap{1}{Kritéria na kontingenčních tabulkách}

\fml{Kontingenční tabulka --- marginály a~očekávané četnosti}
\[
r_k=\sum_{l=1}^{S}a_{kl},\qquad s_l=\sum_{k=1}^{R}a_{kl},\qquad
n=\sum_{k=1}^{R}\sum_{l=1}^{S}a_{kl},\qquad
e_{kl}=\frac{r_k s_l}{n}
\]
\pop{$a_{kl}$ = četnost kombinace $A(v_k) \wedge C(v_l)$ vstupního atributu $A$
($R$ hodnot) a~cílového atributu $C$ ($S$ tříd). Odhady pravděpodobností:
$P(A(v_k)\wedge C(v_l))=a_{kl}/n$, $P(A(v_k))=r_k/n$, $P(C(v_l))=s_l/n$.
$e_{kl}$ je očekávaná četnost \emph{při nezávislosti}.}

\fml{$\chi^2$ jako kritérium relevance a~jako test}
\[
\chi^2=\sum_{k=1}^{R}\sum_{l=1}^{S}\frac{(a_{kl}-e_{kl})^2}{e_{kl}}
   = n\sum_{k=1}^{R}\sum_{l=1}^{S}\frac{\big(a_{kl}-\tfrac{r_ks_l}{n}\big)^2}{r_ks_l}
   = \sum_k\sum_l \frac{a_{kl}^2}{e_{kl}}-n
\]
\pop{\textbf{Čím větší, tím lepší} (relevantnější atribut). Jako test: při
$H_0: P(A=v_k\wedge C=v_l)=P(A=v_k)P(C=v_l)$ má statistika $\chi^2$ rozdělení
s~$(R-1)(S-1)$ stupni volnosti; je-li $\chi^2 \ge \chi^2_{(R-1)(S-1)}(\alpha)$, $H_0$
zamítáme (pro čtyřpolní tabulku a~$\alpha = 0{,}05$ je kritická hodnota $3{,}84$).
Poslední tvar je nejrychlejší pro ruční výpočet. \textbf{Podmínka použitelnosti:}
$e_{kl} = r_ks_l/n \ge 5$ pro všechna $k,l$.}

\fml{Fisherův exaktní test}
\[
p=\frac{r_1!\,r_2!\,s_1!\,s_2!}{n!\,a_{11}!\,a_{12}!\,a_{21}!\,a_{22}!}
     =\frac{\binom{s_1}{a_{11}}\binom{s_2}{a_{12}}}{\binom{n}{r_1}},
\qquad
P=\sum_{i=0}^{m}\frac{r_1!\,r_2!\,s_1!\,s_2!}{n!\,(a_{11}-i)!\,(a_{12}+i)!\,(a_{21}+i)!\,(a_{22}-i)!}
\]
\pop{Pro čtyřpolní tabulky s~\emph{malými} četnostmi, kde $\chi^2$ test nelze použít.
$p$ je pravděpodobnost konkrétní tabulky při daných marginálech; $P$ se sečte přes všechny
uvažované hodnoty ($m=\min(r_1,s_1)$). Je-li $P \le \alpha$, $H_0$ o~nezávislosti se
zamítá.}

\podpodkap{2}{Kritéria založená na entropii}

\fml{Entropie atributu, vzájemná informace, informační míra závislosti}
\[
H(A)=\sum_{k=1}^{R}\frac{r_k}{n}\,H\big(A(v_k)\big),\qquad
H\big(A(v_k)\big)=-\sum_{l=1}^{S}\frac{a_{kl}}{r_k}\log_2\frac{a_{kl}}{r_k}
\]
\[
\mathrm{MI}(A,C)=\sum_{k=1}^{R}\sum_{l=1}^{S}\frac{a_{kl}}{n}\log_2\frac{n\,a_{kl}}{r_k s_l},
\qquad
\mathrm{ID}(A,C)=\frac{\mathrm{MI}(A,C)}{H(C)}
   =\frac{\mathrm{MI}(A,C)}{-\sum_{l}\frac{s_l}{n}\log_2\frac{s_l}{n}}
\]
\pop{$H(A)$ je vážený průměr entropií jednotlivých hodnot --- \textbf{čím menší, tím lepší}
atribut. $\mathrm{MI}$ i~$\mathrm{ID}$ naopak \textbf{čím větší, tím lepší};
$\mathrm{ID}$ je vzájemná informace normalizovaná entropií cílového atributu, takže
je srovnatelná napříč úlohami.}

\podpodkap{3}{Kritéria pro numerické atributy: korelace}

\fml{Pearsonův korelační koeficient}
\[
r=\frac{\mathrm{cov}(X,Y)}{s_X\,s_Y}
   =\frac{\sum_{i=1}^{n}(x_i-\bar x)(y_i-\bar y)}
          {\sqrt{\sum_{i=1}^{n}(x_i-\bar x)^2\ \sum_{i=1}^{n}(y_i-\bar y)^2}}\ \in[-1,1]
\]
\pop{$s_X, s_Y$ jsou směrodatné odchylky, $\bar x, \bar y$ průměry. Míra \textbf{lineárního}
vztahu: $r=\pm1$ = body přesně na přímce, $r=0$ = žádná \emph{lineární} závislost.
$r^2$ (\emph{koeficient determinace}) = podíl rozptylu $Y$ vysvětlený lineární závislostí
na~$X$. Vztah k~regresi: $\beta=\mathrm{cov}(X,Y)/s_X^2$, tedy $r=\beta\,s_X/s_Y$.
\textbf{Tři pasti:} měří jen \emph{lineární} vztah (pro $y=x^2$ symetricky kolem nuly vyjde
$r=0$, přestože je závislost funkční --- nulová korelace \emph{není} nezávislost); je citlivý
na odlehlé hodnoty; a~korelace není kauzalita.}

\fml{Spearmanův koeficient $\rho$ (korelace pořadí)}
\[
\rho = 1-\frac{6\sum_{i=1}^{n} d_i^2}{n\,(n^2-1)}\ \in[-1,1]
\]
\pop{Je to \textbf{Pearson spočítaný z~pořadí} místo z~hodnot; $d_i$ = rozdíl pořadí $i$-tého
objektu v~$X$ a~$Y$. Uvedený zjednodušený tvar platí, \emph{nejsou-li shodná pořadí}
(\emph{ties}). Zachytí \textbf{libovolný monotónní} vztah (ne jen lineární), je robustní vůči
odlehlým hodnotám, \textbf{invariantní vůči monotónní transformaci} ($\log$, $\sqrt{\ }$)
a~funguje i~pro \emph{ordinální} atributy. \emph{Příklad:} $x=(1,2,3,4,5)$, $y=(1,2,4,8,16)$
--- pořadí shodná, všechna $d_i=0$, tedy $\rho=1$, kdežto $r\doteq0{,}93$ (Pearson penalizuje
nelinearitu). Alternativa: Kendallovo $\tau$ přes konkordantní/diskordantní dvojice.}

\podkap{9}{Vztahy mezi atributy: regrese}

\fml{Lineární regrese jedné proměnné (metoda nejmenších čtverců)}
\[
E=\sum_{i=1}^{p}\big(y_i-\beta x_i-\alpha\big)^2
\ \longrightarrow\ \min
\qquad\Longrightarrow\qquad
\beta=\frac{\sum_{i}(x_i-\bar x)(y_i-\bar y)}{\sum_{i}(x_i-\bar x)^2},
\quad \alpha=\bar y-\beta\bar x
\]
\pop{Odvození: položíme $\partial E/\partial\alpha = \partial E/\partial\beta = 0$.
$\beta$ je \emph{výběrová kovariance dělená výběrovým rozptylem $x$}. Predikce
$\hat y = \beta x + \alpha$; regresní přímka vždy prochází bodem $(\bar x, \bar y)$.}

\fml{Vícerozměrná lineární regrese --- normální rovnice}
\[
E=(Y-X\beta)^{\!\top}(Y-X\beta), \qquad
X^{\!\top}\!X\beta=X^{\!\top}Y \qquad\Longrightarrow\qquad
\beta=\big(X^{\!\top}\!X\big)^{-1}X^{\!\top}Y
\]
\pop{$Y$ je vektor cílových hodnot, $X$ rozšířená matice vstupních vzorů (první sloupec jedničky),
$\beta=(\beta_0,\dots,\beta_m)$. Uzavřené řešení, ale invertování $X^{\!\top}X$ je pro
velké úlohy nákladné $\Rightarrow$ iterativní (gradientní) metody.}

\kap{3}{Metody pro dobývání znalostí}

\podkap{1}{Asociační pravidla a~frekventované vzory}

\podpodkap{2}{Support, confidence a~lift}

\fml{Support, confidence, lift}
\begin{gather*}
\mathrm{sup}(R)=\frac{\text{počet transakcí obsahujících } i \text{ a } j}
   {\text{počet všech transakcí}}\cdot 100\,\%,\\
\mathrm{conf}(R)=\frac{\text{počet transakcí obsahujících } i \text{ a } j}
   {\text{počet transakcí obsahujících } i}\cdot 100\,\%,\\
\mathrm{lift}(R)=\frac{p(i\wedge j)}{p(i)\,p(j)}
\end{gather*}
\pop{Pro pravidlo $R: i \Rightarrow j$ nad množinou transakcí.
\emph{Support} --- jak často lze pravidlo použít; \emph{confidence} --- jak spolehlivé jsou
jeho výsledky (odhad $P(j \mid i)$); \emph{lift} --- kolikrát je lepší použít pravidlo než
jen předpokládat závěr. \textbf{Je-li $\mathrm{lift} < 1$, je pravidlo horší než náhodná
volba} a~negace závěru může dát lepší pravidlo; $\mathrm{lift}=1$ znamená nezávislost.}

\podpodkap{4}{Algoritmus Apriori}

\fml{Apriori --- rozklad kombinace na pravidla}
\[
\mathit{Suc}=\mathit{Comb}\setminus \mathit{Ant},\qquad
\mathit{Ant}\cap \mathit{Suc}=\emptyset,\quad \mathit{Ant}\cup \mathit{Suc}=\mathit{Comb}
\]
\pop{Po nalezení frekventovaných kombinací se každá rozdělí na všechny možné dvojice
podkombinací. \emph{Support} pravidla $\mathit{Ant} \Rightarrow \mathit{Suc}$ odpovídá
frekvenci $\mathit{Comb}$, \emph{confidence} poměru frekvencí $\mathit{Comb}$
a~$\mathit{Ant}$. Základ algoritmu je \textbf{uzavřenost dolů} (\emph{downward closure}):
frekvence nadkombinace $\le$ frekvence kombinace, takže nadmnožinu nefrekventované množiny
lze prořezat.}

\podpodkap{5}{Vícenásobný minimální support: MS-Apriori}

\fml{MS-Apriori --- vícenásobný minimální support}
\[
\mathrm{minsup}(\text{pravidla}) = \min_{i \in s} \mathrm{MIS}(i),
\qquad
\max_{i\in s}\mathrm{sup}(i)-\min_{i\in s}\mathrm{sup}(i)\le\varphi
\]
\pop{Řeší \emph{problém vzácných položek} (chléb vs.\ kuchyňský robot): každá položka má
vlastní $\mathrm{MIS}$, minsup pravidla je minimum přes jeho položky, a~rozdíl supportů
uvnitř množiny je omezen $\varphi$. \textbf{Důsledek: uzavřenost dolů neplatí}
--- $\{1,2\}$ může být nefrekventovaná, a~přesto $\{1,2,3\}$ frekventovaná; položky se proto
třídí vzestupně podle MIS a~testuje se proti MIS \emph{první} položky.}

\podkap{2}{Přístupy založené na principu učení s~učitelem}

\podpodkap{2}{Rozhodovací stromy}

\fml{Entropie a~vážená průměrná entropie atributu}
\[
H=-\sum_{j=1}^{m}p_j\log_2 p_j,
\qquad
H\big(A(v)\big)=-\sum_{j=1}^{m}p_j^{A(v)}\log_2 p_j^{A(v)},
\qquad
H(A)=\sum_{v}\frac{|A(v)|}{|D|}\,H\big(A(v)\big)
\]
\pop{$p_j$ = relativní frekvence třídy $j$ na uvažované množině vzorů, $m$ = počet tříd.
Konvence $0\log_2 0 = 0$; přepočet $\log_2 x = \log_{10}x/0{,}301$.
\textbf{Vybírá se atribut s~minimální $H(A)$.} Čistý uzel má $H=0$, dvě stejně velké třídy
dají $H=1$.}

\fml{Informační zisk a~poměrný informační zisk (ID3, C4.5)}
\[
\mathrm{GAIN}(D,A)=H(D)-\sum_{v}\frac{|D_v|}{|D|}H(D_v),
\qquad
\mathrm{GAIN\_RATIO}(D,A)=\frac{\mathrm{GAIN}(D,A)}
   {-\sum_{v}\frac{|D_v|}{|D|}\log_2\frac{|D_v|}{|D|}}
\]
\pop{Očekávané snížení entropie po rozdělení podle atributu; bere se \textbf{největší}.
Protože $H(D)$ je pro všechny atributy stejná, je maximalizace GAIN ekvivalentní minimalizaci
$H(A)$. Jmenovatel GAIN\_RATIO (\emph{split info}) penalizuje \textbf{nadměrné větvení} ---
bez něj by vyhrával atribut jako \uv{rodné číslo}, který dělí data na samé jednoprvkové
podmnožiny.}

\fml{Gini index a~kritérium dělení v~CART}
\[
G=\sum_{i=1}^{r}\frac{n_i}{n}G(v_i),\qquad G(v_i)=1-\sum_{j=1}^{m}p_j^2,
\qquad
\Phi=2P_LP_R\sum_{j=1}^{m}\big|P(C_j\mid t_L)-P(C_j\mid t_R)\big|
\]
\pop{CART generuje \emph{binární} stromy. $n_i$ = počet vzorů s~hodnotou $v_i$,
$p_j$ = pravděpodobnost třídy $j$ v~těchto vzorech; \textbf{nižší $G$ = lepší diskriminace}.
$\Phi$ je kritérium volby dělicího bodu: $P_L, P_R$ jsou pravděpodobnosti levého a~pravého
podstromu; $\Phi$ \textbf{nabývá maxima}, jsou-li všechny vzory dané třídy seskupeny v~témže
podstromu.}

\fml{C4.5 --- pesimistický odhad přesnosti pravidla}
\[
\mathit{acc}-1{,}96\sqrt{\frac{\mathit{acc}(1-\mathit{acc})}{n}},
\qquad \mathit{acc}=\frac{k}{n}
\]
\pop{$k$ = správně klasifikované, $n$ = pokryté vzory. C4.5 hodnotí pravidla \emph{na
trénovacích datech}, kde je přesnost optimistická; charakteristikou pravidla je proto
\emph{dolní mez} intervalu spolehlivosti z~binomického rozdělení (95\% hladina).
Používá se při \emph{rule post-pruning}: vypouštěj předpoklady, dokud se tato mez nezhorší.}

\podpodkap{3}{Bayesovská klasifikace}

\fml{Bayesova formule a~hypotéza MAP}
\[
P(H\mid E)=\frac{P(E\mid H)\,P(H)}{P(E)},
\qquad
H_{MAP}=\arg\max_{H\in\mathcal H} P(H\mid E)=\arg\max_{H\in\mathcal H} P(E\mid H)\,P(H)
\]
\pop{$P(H)$ = apriorní pravděpodobnost hypotézy, $P(H\mid E)$ = aposteriorní, $P(E\mid H)$ =
věrohodnost. Jmenovatel $P(E)$ je pro všechny hypotézy stejný, takže jej lze při
maximalizaci \emph{zanedbat}.}

\fml{Naivní bayesovský klasifikátor}
\[
P(H\mid E_1,\dots,E_k)=\frac{P(H)}{P(E_1,\dots,E_k)}\prod_{i=1}^{k}P(E_i\mid H),
\qquad
H_{MAP}=\arg\max_{H_t} P(H_t)\prod_{i=1}^{k}P(E_i\mid H_t)
\]
\pop{\textbf{Předpoklad:} jevy $E_1,\dots,E_k$ jsou \emph{podmíněně nezávislé} při platnosti
$H$ --- v~reálných úlohách splněn jen zřídka, odtud \uv{naivní}. Odhady z~dat:
$P(H_t) = \frac{\text{počet vzorů z~}t}{\text{počet všech vzorů}}$,
$P(E_k \mid H_t) = \frac{\text{počet vzorů z~}t\text{ splňujících }E_k}{\text{počet vzorů z~}t}$.
Přes nesplněný předpoklad funguje překvapivě dobře, protože pro \emph{argmax} stačí správné
pořadí, ne přesné pravděpodobnosti.}

\podpodkap{4}{Support Vector Machines}

\fml{Lineární SVM --- rozhodovací funkce a~okraj}
\[
f(x)=\langle w\cdot x\rangle+b, \qquad
d_+=d_-=\frac{1}{\|w\|},\qquad \text{margin}=d_++d_-=\frac{2}{\|w\|}
\]
\pop{Trénovací množina $D=\{(x_i,y_i)\}$, $y_i \in \{1,-1\}$; $w$ je normálový vektor a~$b$ posun hyperroviny. Vzdálenost bodu od
hyperroviny je $|\langle w\cdot x_i\rangle+b| / \|w\|$. Separujících hyperrovin je mnoho ---
SVM hledá tu s~\textbf{největším okrajem}, protože ta podle teorie učení minimalizuje mez
chyby.}

\fml{Primární optimalizační úloha a~Lagrangián}
\[
\min\ \frac{\langle w\cdot w\rangle}{2}
   \quad\text{za podmínek}\quad y_i\big(\langle w\cdot x_i\rangle+b\big)\ge 1
\]
\[
L_P=\frac{\langle w\cdot w\rangle}{2}-\sum_{i=1}^{r}\alpha_i
   \big[y_i(\langle w\cdot x_i\rangle+b)-1\big],\qquad \alpha_i\ge 0
\]
\pop{Maximalizace $2/\|w\|$ je ekvivalentní minimalizaci $\langle w\cdot w\rangle/2$.
Podmínka je souhrnný zápis $\langle w\cdot x_i\rangle+b\ge 1$ pro $y_i=1$ a~$\le -1$ pro
$y_i=-1$. Řešení musí splňovat \emph{Kuhn--Tuckerovy podmínky}, které jsou pro konvexní
účelovou funkci s~lineárními omezeními nejen nutné, ale i~\emph{postačující}.}

\fml{Kuhn--Tuckerovy podmínky a~vznik duálu}
\[
\frac{\partial L_P}{\partial w_j}=0 \Rightarrow w=\sum_{i=1}^{r} y_i\alpha_i x_i,
\qquad
\frac{\partial L_P}{\partial b}=0 \Rightarrow \sum_{i=1}^{r} y_i\alpha_i=0,
\qquad
\alpha_i\big(y_i(\langle w\cdot x_i\rangle+b)-1\big)=0
\]
\pop{\textbf{Postup: (1)} vynuluj parciální derivace $L_P$ podle \emph{primárních} proměnných
$w$ a~$b$; \textbf{(2)} dosaď získané vztahy zpět do $L_P$ --- $w$ a~$b$ se vyeliminují
a~zbudou jen $\alpha_i$. KKT podmínky jsou obecně jen nutné, pro konvexní účelovou funkci
s~lineárními omezeními i~\emph{postačující}. Poslední je \textbf{komplementarita}: součin je
nulový, takže je-li $\alpha_i>0$, musí být $y_i(\langle w\cdot x_i\rangle+b)=1$ --- bod leží
\emph{přesně na} okrajové hyperrovině. To jsou \textbf{podpůrné vektory}; pro ostatní je
$\alpha_i=0$ a~řešení neovlivňují.}

\fml{Wolfeho duální úloha (\emph{vázaná}) a~rekonstrukce $w$, $b$}
\[
\text{max } L_D=\sum_{i=1}^{r}\alpha_i-\frac12\sum_{i,j}y_iy_j\alpha_i\alpha_j
   \langle x_i\cdot x_j\rangle
\qquad\text{za podmínek}\qquad
\sum_{i=1}^{r} y_i\alpha_i=0,\ \ \alpha_i\ge 0
\]
\[
w=\sum_{i\in \mathit{sv}} y_i\alpha_i x_i,
\qquad
b = y_k - \langle w\cdot x_k\rangle \ \ (x_k \text{ podpůrný vektor})
\]
\pop{\textbf{Omezení jsou součástí úlohy} --- vznikla v~kroku (1), bez nich je duál nedopsaný.
Maximum $L_D$ se nabývá pro tytéž $w,b,\alpha_i$ jako minimum $L_P$; řeší se numericky
(prakticky SMO). Duál obsahuje vstupy \emph{jen} přes skalární součiny $\langle x_i\cdot
x_j\rangle$ --- právě to umožňuje kernelový trik. Vztah pro $b$ plyne z~komplementarity
($y_k(\langle w\cdot x_k\rangle+b)=1$ a~$y_k^2=1$); numericky se průměruje přes všechny
podpůrné vektory.}

\fml{Rozhodovací hranice a~testování}
\[
\langle w\cdot x\rangle+b=\sum_{i\in \mathit{sv}}y_i\alpha_i\langle x_i\cdot x\rangle+b=0,
\qquad
\text{test: } \mathrm{sign}\Big(\sum_{i\in \mathit{sv}}\alpha_iy_i\langle x_i\cdot z\rangle+b\Big)
\]
\pop{$z$ je testovaný vzor, $\mathit{sv}$ množina podpůrných vektorů. Vrátí-li výraz 1, je $z$
pozitivní, jinak negativní. Všimni si, že $w$ ani $\phi$ se nikde nevyčíslují --- stačí
skalární součiny s~podpůrnými vektory.}

\fml{Kernelový trik a~měkký okraj}
\[
K(x,z)=\langle\phi(x)\cdot\phi(z)\rangle,
\qquad
\min\ \frac{\langle w\cdot w\rangle}{2}+C\sum_i\xi_i, \quad \xi_i\ge 0
\]
\pop{Kernel nahradí skalární součin ve všech výrazech, aniž bychom $\phi$ kdy vyčíslili ---
odstraňuje prokletí dimenzionality při explicitní transformaci. Příklad
($d=2$, 2-D vstup): $\langle x\cdot z\rangle^2 = \big\langle (x_1^2,x_2^2,\sqrt2 x_1x_2)
\cdot(z_1^2,z_2^2,\sqrt2 z_1z_2)\big\rangle$. Zda je funkce kernelem, říká
\emph{Mercerova věta}; běžné kernely: lineární, polynomiální, gaussovský (RBF), sigmoidální.
\emph{Slack variables} $\xi_i$ připouštějí chyby; $C$ řídí kompromis mezi šířkou okraje
a~počtem chyb.}

\podpodkap{5}{Logistická regrese}

\fml{Logit a~sigmoida}
\[
\ln\frac{P(y=1\mid x)}{1-P(y=1\mid x)}=\beta_0+\beta_1x_1+\cdots+\beta_mx_m=z,
\qquad
P(y=1\mid x)=\frac{1}{1+e^{-z}}=\sigma(z)
\]
\pop{$\beta_j$ jsou koeficienty modelu a~$z$ jejich lineární kombinace (\emph{logit}). Modeluje se logaritmus \emph{podmíněné šance} (\emph{conditional odds}) jako lineární
kombinace vstupů. Logistická regrese je nejdůležitější případ nelineární regrese a~zároveň
klasifikátor; výstupem je \emph{pravděpodobnost}, ne jen třída.}

\fml{Učení logistické regrese --- věrohodnost a~křížová entropie}
\[
L=\prod_{i=1}^{p}\Big(\frac{1}{1+e^{-z_i}}\Big)^{y_i}\Big(\frac{1}{1+e^{z_i}}\Big)^{1-y_i},
\qquad
E=-\log L=-\sum_i\Big[y_i\log\sigma(z_i)+(1-y_i)\log\big(1-\sigma(z_i)\big)\Big]
\]
\pop{Místo maximalizace věrohodnosti $L$ se minimalizuje \emph{křížová entropie} $E$
(numericky stabilnější, součet místo součinu).}

\fml{Adaptační pravidla logistické regrese}
\[
\Delta\beta_j=-\eta\frac{\partial E}{\partial\beta_j}
   =\eta\sum_i\big(y_i-\sigma(z_i)\big)x_{i,j},\qquad
\Delta\beta_0=\eta\sum_i\big(y_i-\sigma(z_i)\big)
\]
\pop{$\eta$ je rychlost učení. Odvození využívá $\sigma'=\sigma(1-\sigma)$, po zjednodušení se derivace sigmoidy
\emph{vykrátí}. Chyba tedy vstupuje do úpravy vah přímo jako rozdíl \emph{pozorované}
a~\emph{predikované} hodnoty --- tentýž tvar jako delta pravidlo perceptronu.}

\podpodkap{6}{Diskriminační analýza}

\fml{Diskriminační analýza --- diskriminační funkce}
\[
x \in C_j \iff g_j(x)=\max_{k} g_k(x),
\qquad
g_j(x)=\beta_{j0}+\beta_{j1}x_1+\cdots+\beta_{jm}x_m \ \text{(lineární DA)}
\]
\pop{Optimální klasifikace ve smyslu minimální chyby nastane, jsou-li diskriminační funkce
úměrné aposteriorním pravděpodobnostem: $g_j(x)=P(C_j\mid x)=p(x\mid C_j)P(C_j)/p(x)$;
pro dvě třídy $g(x)=p(x\mid C_1)P(C_1)-p(x\mid C_2)P(C_2)$.}

\fml{LDA vs.\ QDA pro normální rozdělení}
\[
g(x)=(\mu_1-\mu_2)^{\!\top}S^{-1}x-\tfrac12(\mu_1-\mu_2)^{\!\top}S^{-1}(\mu_1+\mu_2)
   -\ln\frac{P(C_2)}{P(C_1)}
\]
\[
g_j(x)=\ln p_j-\tfrac12\,(x-\mu_j)^{\!\top}\Sigma^{-1}(x-\mu_j)
\]
\pop{Pro normální rozdělení s~\emph{obecnými} kovariančními maticemi $S_1 \ne S_2$ vychází
\textbf{kvadratická} diskriminační funkce (QDA); jsou-li matice \emph{shodné} ($S_1=S_2=S$),
redukuje se na \textbf{lineární} (LDA) --- první vzorec. Druhý je ekvivalentní tvar LDA se
společnou $\Sigma$; odhaduje se jen $\mu_j$ ($q$ parametrů na třídu), a~hranice mezi dvěma
třídami $\{x: g_j(x)=g_k(x)\}$ je \emph{nadrovina}. Jsou-li navíc třídy stejně pravděpodobné
a~$S=I$, jde o~prosté srovnávání vzdáleností od středních hodnot.}

\podpodkap{7}{Neuronové sítě a~extrémní učicí stroje (ELM)}

\fml{ELM --- extrémní učicí stroj}
\[
f_L(x)=\sum_{i=1}^{L}\beta_i\,G(a_i,b_i,x),
\qquad
H\beta=T \qquad\Longrightarrow\qquad \beta = H^{+}T
\]
\pop{Dopředná síť s~\emph{jednou} skrytou vrstvou o~$L$ neuronech, $\beta_i$ jsou výstupní váhy; $G$ je např.\ sigmoida
$g(\langle a_i\cdot x\rangle+b_i)$ nebo RBF $g(b_i\|x-a_i\|)$.
\textbf{Klíčová myšlenka:} parametry skrytých neuronů $(a_i,b_i)$ se \emph{neadaptují} ---
vygenerují se \emph{náhodně}, i~bez znalosti dat (s~pravděpodobností 1 platí
$\lim_{L\to\infty}\|f-f_L\|=0$). Učí se jen výstupní váhy, a~to \emph{jednou},
Mooreovou--Penroseovou pseudoinverzí $H^{+}$ výstupní matice skryté vrstvy $H$.
Neiterativní $\Rightarrow$ rychlé, bez lokálních minim; regularizace přičtením $I/C$
na diagonálu $H^{\!\top}H$.}

\podpodkap{8}{Ensemble metody}

\fml{Bagging --- pokles rozptylu a~bootstrap}
\[
\mathrm{var}\Big(\tfrac1k\sum_{i=1}^k h_i\Big) = \frac{\sigma^2}{k} \quad\text{(i.i.d.)},
\qquad
1-\Big(1-\frac1n\Big)^{\!n}\approx 1-\frac1e=63{,}2\,\%
\]
\pop{Bagging snižuje \emph{rozptyl}. Bootstrapový vzorek (výběr $n$ bodů \emph{s~opakováním})
obsahuje asi 63,2~\% různých bodů originálu --- pravděpodobnost, že bod ve vzorku \emph{není},
je $(1-1/n)^n \to 1/e$. Ideální bázový model: rozhodovací stromy (nízké zkreslení, vysoký
rozptyl). \textbf{Omezení: bagging nesnižuje zkreslení.}}

\fml{Náhodné lesy --- rozptyl při korelovaných komponentách}
\[
\mathrm{var} = \rho\sigma^2+\frac{(1-\rho)\sigma^2}{k}
\]
\pop{$\rho$ je vzájemná pozitivní korelace $k$ klasifikátorů. Člen $\rho\sigma^2$ je
\textbf{nezávislý na $k$} --- právě on omezuje přínos baggingu, protože bootstrapované stromy
bývají korelované. Náhodný les korelaci snižuje tím, že v~každém uzlu předloží jen
\emph{náhodnou podmnožinu} $m \approx \sqrt{\text{počet příznaků}}$ příznaků.}

\fml{AdaBoost --- váha komponenty a~převáženi vzorů}
\[
\alpha_t = \tfrac12\ln\frac{1-\epsilon_t}{\epsilon_t},
\qquad
W_{t+1}(i) = \frac{W_t(i)\,e^{\pm\alpha_t}}{\sum_j W_{t+1}(j)}
\]
\pop{$\epsilon_t$ = vážená chyba v~$t$-té iteraci, $W_1(i)=1/n$. Znaménko v~exponentu:
$+\alpha_t$ pro \emph{nesprávně} klasifikované (jejich váha roste), $-\alpha_t$ pro správně
klasifikované. Konec, když $t \ge T$, $\epsilon_t = 0$ nebo $\epsilon_t \ge 0{,}5$.
\textbf{Boosting snižuje zkreslení} (bagging rozptyl), musí ale běžet \emph{sekvenčně}
a~je citlivý na šum.}

\podkap{3}{Klastrová analýza}

\podpodkap{1}{Formulace úlohy a~míry vzdálenosti}

\fml{Metriky vzdálenosti}
\[
d_H=\sum_{i=1}^{m}|x_{1i}-x_{2i}|,\quad
d_E=\sqrt{\sum_{i=1}^{m}(x_{1i}-x_{2i})^2},\quad
d_C=\max_i |x_{1i}-x_{2i}|,\quad
d_q=\Big(\sum_{i=1}^{m}|x_{1i}-x_{2i}|^{q}\Big)^{1/q}
\]
\pop{Hammingova (Manhattan, $L_1$), euklidovská ($L_2$), Čebyševova ($L_\infty$)
a~Minkowského ($L_q$), jejímiž speciálními případy jsou první tři:
$d_H=d_1$, $d_E=d_2$, $d_C=\lim_{q\to\infty}d_q$. Všechny předpokládají \emph{stejný
rozptyl} veličin $\Rightarrow$ nutná normalizace.}

\fml{Mahalanobisova vzdálenost}
\[
d_M^2(x_1,x_2)=(x_1-x_2)^{\!\top}S^{-1}(x_1-x_2)
\]
\pop{$S$ je kovarianční matice. Zohledňuje \emph{různé rozptyly i~korelace} mezi veličinami;
pro $S=I$ se redukuje na kvadrát euklidovské vzdálenosti. Používá se, mají-li veličiny různý
rozptyl.}

\fml{Vzdálenost mezi klastry (metody spojování)}
\[
D_{\min}=\min_{x_i,x_j} d(x_i,x_j),\quad
D_{\max}=\max_{x_i,x_j} d(x_i,x_j),\quad
D_{\mathrm{avg}}=\frac{1}{n_kn_l}\sum_{i=1}^{n_k}\sum_{j=1}^{n_l} d(x_i,x_j)
\]
\pop{\emph{Single linkage} (nejbližší soused --- řetězí, umí protáhlé tvary),
\emph{complete linkage} (nejvzdálenější soused --- kompaktní kulovité klastry),
\emph{average linkage} (kompromis), \emph{centroidní} $D=d(c_k,c_l)$.
Volba metody určuje tvar nalezených klastrů v~hierarchickém shlukování.}

\podpodkap{2}{Metoda $k$-means}

\fml{$k$-means --- centroid a~$k$-medoids --- účelová funkce}
\[
c_k=\frac{1}{|C_k|}\sum_{x\in C_k}x,
\qquad
O=\sum_{i=1}^{n}\min_{j} d(x_i,y_j)
\]
\pop{$k$-means iteruje \uv{spočti centroidy $\to$ přiřaď každý vzor k~nejbližšímu}, dokud
dochází k~přesunům. $k$-medoids hledá $k$ reprezentantů $y_1,\dots,y_k$ \emph{z~původní
databáze} minimalizujících $O$ --- je proto robustní vůči odlehlým hodnotám a~funguje pro
libovolný typ dat, ale je pomalejší; složitost jedné iterace $O(k\,n\,d)$.
CLARA/CLARANS jej škálují vzorkováním, resp.\ náhodným prohledáváním.}

\podpodkap{5}{Vektorová kvantizace: algoritmus LVQ}

\fml{LVQ --- adaptace vah (shlukování s~učitelem)}
\[
q = \arg\min_j \|x-w_j\|_2^2,
\qquad
w_q(k+1)= w_q(k) \pm \mu(k)\big(x-w_q(k)\big)
\]
\pop{Znaménko $+$ (\uv{přitáhni}), souhlasí-li třída $x$ s~třídou vítězného vektoru $w_q$,
jinak $-$ (\uv{odtlač}). Rychlost učení klesá, např.\ $\mu(k)=\mu(k-1)/(k+1)$.
Váhové vektory se inicializují prvními $m$ vzory (musí pokrýt všechny třídy), čímž se
neurony \uv{zkalibrují} s~třídami.}

\kap{4}{Charakteristická a~diskriminační pravidla a~jejich zajímavost}

\podkap{1}{Pravidla jako srozumitelná forma klasifikace}

\fml{Fuzzy implikace jako reziduum t-normy}
\[
t(\varphi\Rightarrow\psi)=\sup\{\,c\in[0,1]: T(t(\varphi),c)\le t(\psi)\,\},
\qquad
t(\varphi\Leftrightarrow\psi)=T\big(t(\varphi\Rightarrow\psi),\,t(\psi\Rightarrow\varphi)\big)
\]
\pop{Booleovská pravidla pracují s~$t(\varphi)\in\{0,1\}$, fuzzy pravidla se stupni
pravdivosti $t(\varphi)\in[0,1]$. Na objektech, kde antecedent \emph{plně} platí
($t(\varphi)=1$), implikace splývá s~ekvivalencí --- proto se v~praxi pravidla vyhodnocují
právě tam.}

\podkap{3}{Charakteristická a~diskriminační pravidla: t-weight a~d-weight}

\fml{t-weight (typičnost) a~kvantitativní charakteristické pravidlo}
\[
t\text{-}\mathrm{weight}(q_a)=\frac{\mathrm{count}(q_a)}{\sum_{i=1}^{N}\mathrm{count}(q_i)}\in[0,1]
\]
\[
\forall X\ \ \mathrm{target}(X)\Rightarrow \mathrm{cond}_1(X)\,[t{:}\,w_1]\ \vee\
   \cdots\ \vee\ \mathrm{cond}_N(X)\,[t{:}\,w_N]
\]
\pop{$q_1,\dots,q_N$ jsou zobecněné záznamy cílové třídy, $X$ objekt (pravidlo platí pro všechny). t-weight = \emph{podíl objektů
cílové třídy, které $q_a$ pokrývá} (součet přes všechny záznamy je 100~\%). Charakteristické
pravidlo je \textbf{nutná podmínka} --- implikace \emph{z}~třídy. Odpovídá \emph{recallu}
pravidla vůči cílové třídě.}

\fml{d-weight (rozlišovací váha) a~kvantitativní diskriminační pravidlo}
\[
d\text{-}\mathrm{weight}(q_a)=\frac{\mathrm{count}(q_a\in C_t)}
   {\mathrm{count}(q_a\in C_t)+\sum_{j=1}^{m}\mathrm{count}(q_a\in C_j)}\in[0,1]
\]
\[
\forall X\ \ \mathrm{target}(X)\Leftarrow \mathrm{cond}(X)\,[d{:}\,w]
\]
\pop{$C_t$ je cílová třída, $C_1,\dots,C_m$ kontrastní třídy. d-weight = \emph{podíl objektů
pokrytých záznamem $q_a$, které patří do cílové třídy}. Diskriminační pravidlo je
\textbf{postačující podmínka} --- implikace \emph{do}~třídy. Odpovídá \emph{precision}
pravidla. \emph{Kvantitativní popisné pravidlo} kombinuje obojí:
$\Leftrightarrow$ s~$[t{:}\,w,\ d{:}\,w']$.}

\podkap{6}{Observační pravidla: konfidence, podpora a~testy hypotéz}

\fml{Observační pravidlo s~konfidencí a~podporou}
\[
\mathrm{conf}=\frac{a}{a+b},\qquad \mathrm{sup}=\frac{a}{n},\qquad
(\varphi\sim\psi)=1 \iff \mathrm{conf}\ge c_0\ \wedge\ \mathrm{sup}\ge s_0
\]
\pop{Čtyřpolní tabulka: $a$ objektů splňuje $\varphi\wedge\psi$, $b$ objektů
$\varphi\wedge\neg\psi$, $c$ objektů $\neg\varphi\wedge\psi$, $d$ objektů
$\neg\varphi\wedge\neg\psi$, $n=a+b+c+d$. $\mathrm{conf}$ je nestranný odhad
$P(\psi\mid\varphi)$. Druhý typ observačních pravidel je založen na \emph{testech hypotéz}
místo prahů. Každá volba míry a~prahu je jiný \emph{zobecněný kvantifikátor}, a~tedy jiná
množina nalezených pravidel.}

\podkap{8}{Měření zajímavosti pravidel}

\fml{Střední hodnota ceny (různá cena různých chyb)}
\[
\frac1n\sum_{i,j}c(i,j)\,n_{ij}
\]
\pop{$c(i,j)$ = cena klasifikace třídy $i$ jako $j$, $n_{ij}$ = odpovídající četnost.
Tradiční chybovost je speciální případ: $c(i,i)=0$, jinak 1. Realističtější, protože
nedoručený ham vadí víc než nerozpoznaný spam a~nerozpoznané napadení sítě víc než falešný
poplach.}

\kap{5}{Reprezentace, vyhodnocování a~vizualizace znalostí}

\podkap{2}{Vyhodnocování modelů}

\podpodkap{1}{Kritéria hodnocení klasifikačních metod}

\fml{Přesnost a~křížová validace}
\[
\mathit{Accuracy}=\frac{\text{počet správných klasifikací}}{\text{počet všech testovaných případů}},
\qquad
\mathit{Accuracy}_{CV}=\frac1k\sum_{i=1}^{k}\mathit{Accuracy}_i
\]
\pop{\textbf{Zásadní pravidlo:} trénovací množina se nesmí použít při testování a~testovací
při učení --- jen neviděná data dají nezkreslený odhad. \emph{Holdout} pro velká data,
\emph{$k$-násobná křížová validace} ($k=5$ nebo 10) pro menší, \emph{leave-one-out} pro velmi
malá. \emph{Validační} množina (třetí část) slouží k~ladění \emph{parametrů} algoritmu.}

\fml{Interval spolehlivosti pro křížovou validaci}
\[
\mathit{Accuracy}_{CV}\pm t_{N,k-1}\cdot s_{\mathit{Accuracy}},
\qquad
s_{\mathit{Accuracy}}=\sqrt{\frac{1}{k(k-1)}\sum_{i=1}^{k}
   \big(\mathit{Accuracy}_i-\mathit{Accuracy}_{CV}\big)^2}
\]
\pop{Přibližný $N\%$ interval; $t_{N,k-1}$ je kvantil Studentova rozdělení
s~$k-1$ stupni volnosti. Umožňuje říct, zda je rozdíl mezi dvěma modely
\emph{statisticky významný}, nebo jen šum.}

\podpodkap{3}{Matice záměn a~odvozené míry\label{ssec:matice-zamen}}

\fml{Matice záměn: precision, recall, $F_1$}
\[
p=\mathit{precision}=\frac{TP}{TP+FP},\qquad
r=\mathit{recall}=\frac{TP}{TP+FN},\qquad
F_1=\frac{2pr}{p+r}=\frac{2}{\frac1p+\frac1r}
\]
\pop{$TP$/$TN$ = správně klasifikované pozitivní/negativní vzory, $FP$/$FN$ = nesprávně.
\emph{Precision} = podíl správných mezi vzory \emph{označenými} za pozitivní;
\emph{recall} = podíl nalezených mezi \emph{skutečně} pozitivními. \textbf{Obě hodnotí jen
pozitivní třídu} --- proto se používají u~nevyvážených dat, kde je accuracy zavádějící.
$F_1$ je \emph{harmonický průměr}, a~je tedy blíž tomu \emph{menšímu} z~obou čísel: velký
je jen tehdy, jsou-li velké \emph{obě} míry.}

\podpodkap{5}{ROC křivka a~AUC}

\fml{ROC křivka: TPR, FPR, TNR}
\[
\mathit{TPR}=\frac{TP}{TP+FN},\qquad \mathit{FPR}=\frac{FP}{FP+TN},\qquad
\mathit{TNR}=\frac{TN}{FP+TN}
\]
\pop{ROC je graf TPR proti FPR při měnícím se prahu skóre. $\mathit{TPR}$ je \emph{recall}
pozitivní třídy (též \emph{senzitivita}), $\mathit{TNR}$ \emph{specificita}
($=1-\mathit{FPR}$). \textbf{AUC} = plocha pod křivkou: 0,5 = náhodný klasifikátor,
1,0 = dokonalý; interpretuje se jako pravděpodobnost, že náhodný pozitivní vzor dostane vyšší
skóre než náhodný negativní. Na rozdíl od accuracy nezávisí na zvoleném prahu ani na poměru
tříd.}

\kap{6}{Modely pro analýzu sociálních sítí, centralita a~detekce komunit}

\podkap{3}{Síla vazeb}

\fml{Překryv sousedství (síla vazby)}
\[
O(i,j)=\frac{|N_i\cap N_j|}{|N_i\cup N_j|-2}
\]
\pop{$N_i$ = sousedé uzlu $i$; $-2$ ve jmenovateli vylučuje $i$ a~$j$ samotné.
\textbf{Čím větší překryv, tím silnější vazba.} Souvisí se \emph{sílou slabých vazeb}
(Granovetter): mosty a~zkratkové mosty mají překryv blízký nule, ale zprostředkovávají
nové informace. \emph{Zkratkový most}: hrana, jejímž odstraněním se vzdálenost koncových
uzlů zvětší na více než 2 --- čím větší výsledná vzdálenost, tím slabší vazba.}

\podkap{4}{Klíčoví hráči: míry centrality}

\fml{Centrality}
\[
C_D(v) = \deg(v), \qquad
C_B(v) = \sum_{s \ne v \ne t} \frac{\sigma_{st}(v)}{\sigma_{st}}, \qquad
C_C(v) = \frac{n-1}{\sum_{u} d(v,u)}
\]
\pop{\emph{Stupňová} ($C_D$) --- kolik lidí osoba dosáhne \textbf{přímo}; propojenost,
vliv, popularita. \emph{Mezilehlá} ($C_B$) --- $\sigma_{st}$ je počet nejkratších cest mezi
$s$ a~$t$, $\sigma_{st}(v)$ počet těch přes $v$; indikuje uzly na komunikačních cestách,
tedy místa, kde se síť může \textbf{rozpadnout}. \emph{Blízkostní} ($C_C$, kde $d(v,u)$ je vzdálenost uzlů a~$n$ počet uzlů sítě) --- převrácená
průměrná délka nejkratších cest; míra \emph{dosahu} a~rychlosti šíření.
\textbf{Jednotlivé míry nestačí:} dva uzly \emph{společně} mohou být pro síť důležitější než
nejcentrálnější uzel sám.}

\fml{Vlastní centralita (\emph{eigenvector centrality})}
\[
x_i=\frac1\lambda\sum_{j\in M(i)}x_j=\frac1\lambda\sum_{j\in V}a_{i,j}x_j
   \quad\Longleftrightarrow\quad Ax=\lambda x
\]
\pop{$A$ je matice sousednosti, $M(i)$ množina sousedů $i$, $x_i$ centralita uzlu $i$. Centralita uzlu je úměrná součtu centralit jeho sousedů --- moc a~status jsou definovány
\emph{rekurzivně}. Vlastních čísel je mnoho, ale \emph{požadavek nezápornosti} všech složek
vlastního vektoru vybírá právě \textbf{největší} $\lambda$ (Perronova--Frobeniova věta).
Vektor je určen až na násobek $\Rightarrow$ nutná normalizace. Blízce souvisí s~PageRankem.}

\podkap{5}{Koheze: míry na úrovni sítě}

\fml{Reciprocita, hustota, klastrovací koeficient}
\[
\text{density}=\frac{|E|}{|V|(|V|-1)/2} \ \text{(neor.)},\quad
\frac{|E|}{|V|(|V|-1)} \ \text{(or.)},
\qquad
C_i=\frac{2\,|\{e_{jk}: v_j,v_k\in N_i\}|}{k_i(k_i-1)}
\]
\pop{\emph{Reciprocita} = podíl obousměrných vztahů; \textbf{má smysl jen v~orientovaných
grafech}. \emph{Hustota} = podíl hran na počtu \emph{možných}; klika má hustotu 1.
\emph{Klastrovací koeficient} uzlu = hustota jeho \emph{sousedství} $N_i$, $k_i=|N_i|$
(u~orientovaných grafů bez faktoru 2); koeficient sítě $C=\frac{1}{|V|}\sum_i C_i$
indikuje přítomnost (sub)komunit.}

\podkap{7}{Modelování vlivu}

\fml{Model lineárního prahu (LTM)}
\[
\sum_{u\in N_v^{\text{aktivní}}} w_{u,v}\ \ge\ \theta_v \sum_{u\in N_v} w_{u,v}
\]
\pop{$w_{u,v}$ je váha vazby ze souseda $u$ na $v$. Každý uzel $v$ si zvolí prah $\theta_v$ \emph{náhodně} z~$U[0,1]$; aktivní uzly zůstávají
aktivní a~v~každém kroku se aktivují uzly splňující tuto podmínku. LTM se soustředí na
\textbf{příjemce} vlivu (ten sčítá tlak okolí); nezávislý kaskádový model (ICM) naopak na
\emph{odesílatele} --- každý nově aktivovaný uzel dostane jediný pokus aktivovat každého
souseda.}

\fml{Maximalizace vlivu}
\[
\max_{|B|\le k}\ \sigma(B)
\]
\pop{$\sigma(B)$ = očekávaný počet uzlů ovlivněných počáteční aktivační množinou $B$.
Úloha je \textbf{NP-těžká}, ale $\sigma$ je submodulární, takže \emph{hladový} přístup
(přidávej uzel s~největším marginálním přírůstkem) dává řešení na úrovni
\textbf{63~\%} ($1-1/e$) optima.}

\podkap{8}{Vliv versus korelace}

\fml{Logistický model vlivu}
\[
p(a)=\frac{e^{\alpha\ln(a+1)+\beta}}{1+e^{\alpha\ln(a+1)+\beta}},\qquad
\ln\frac{p(a)}{1-p(a)}=\alpha\ln(a+1)+\beta
\]
\pop{$a$ = počet aktivních kontaktů uzlu, $\alpha$ = \emph{koeficient sociální korelace},
$\beta$ = vrozená tendence k~aktivaci, $\ln(a+1)$ = vysvětlující proměnná (logaritmus tlumí
vliv velkého počtu kontaktů). Jde o~logistickou regresi; \emph{shuffle test} (přeházení
časových značek) rozliší skutečný vliv od pouhé korelace.}

\podkap{9}{Homofilie, výběr a~sociální vliv}

\fml{Test homofilie}
\[
\text{očekávaný podíl heterogenních hran} = 2pq
\]
\pop{$p$, $q$ jsou podíly obou hodnot atributu v~síti; při náhodném přiřazení je
pravděpodobnost hrany muž--muž $p^2$, žena--žena $q^2$, různopohlavní $2pq$.
\textbf{Je-li podíl heterogenních hran významně \emph{nižší} než $2pq$, jde o~důkaz
homofilie}; významně \emph{vyšší} znamená \emph{inverzní} homofilii (partnerské vztahy).
Příklad: 5 z~18 hran různopohlavních, $p=2/3$, $q=1/3$, $2pq=4/9$, tj.\ očekáváno 8 z~18
$\Rightarrow$ homofilie. Test funguje i~pro více než dvě hodnoty.}

\podkap{13}{Analýza vazeb: huby, autority a~algoritmus HITS}

\fml{HITS --- huby a~autority}
\[
a(i)=\sum_{(j,i)\in E}h(j),\qquad h(i)=\sum_{(i,j)\in E}a(j),
\qquad\text{maticově}\qquad
a=L^{\!\top}h,\qquad h=La
\]
\pop{$L$ je matice sousednosti hypertextového grafu \emph{bázové množiny}
(kořenová množina $S$ z~vyhledávače $+$ stránky odkazující na $S$ a~odkazované z~$S$).
$a(i)$ je skóre autority a~$h(i)$ skóre hubu stránky $i$; vzájemné zesilování: dobrá autorita je odkazovaná dobrými huby a~naopak.
\textbf{Normalizace je nezbytná} --- bez ní hodnoty rostou nade všechny meze
a~mocninná iterace nekonverguje.}

\fml{HITS --- spektrální analýza}
\[
a_k=\big(L^{\!\top}L\big)^{k-1}a_0,\qquad h_k=\big(LL^{\!\top}\big)^{k}h_0,
\qquad
\frac{a_k}{c_1^{k}}=x_1z_1+x_2\Big(\frac{c_2}{c_1}\Big)^{\!k}z_2+\cdots
\]
\pop{$L^{\!\top}L$ je \emph{symetrická}, má tedy $n$ vlastních vektorů $z_1,\dots,z_n$
s~vlastními čísly $|c_1|\ge\cdots\ge|c_n|$; rozvineme-li $a_0=\sum_i x_iz_i$, jde pro
$c_1 > c_2$ a~$k\to\infty$ každý člen kromě prvního k~nule. Vektor autorit tedy konverguje
k~\textbf{hlavnímu vlastnímu vektoru} $L^{\!\top}L$ --- rychlost konvergence určuje poměr
$c_2/c_1$.}

\podkap{14}{PageRank}

\fml{Prestiž pořadí a~základní PageRank}
\[
P(i)=A_{1i}P(1)+\cdots+A_{ni}P(n) \iff P=A^{\!\top}P,
\qquad
P(i)=\sum_{(j,i)\in E}\frac{P(j)}{O_j}
\]
\pop{Prestiž je lineární kombinace pořadí těch, kteří na $i$ odkazují --- \emph{charakteristická
rovnice} vlastního systému matice $A^{\!\top}$. PageRank navíc \textbf{dělí} skóre stránky
mezi její odkazy: $O_j$ = počet výstupních odkazů $j$, tedy $A_{ij}=1/O_i$ pro $(i,j)\in E$.
Hlas důležité stránky váží víc, ale stránka odkazující na tisíc jiných dá každé málo.}

\fml{Finální PageRank s~tlumicím faktorem}
\[
P=\Big(\frac{1-d}{n}\Big)\,\mathbf{E}+d\,A^{\!\top}P
\qquad\Longrightarrow\qquad
P(i)=(1-d)+d\sum_{(j,i)\in E}\frac{P(j)}{O_j}
\]
\pop{$\mathbf E = \mathbf{ee}^{\!\top}$ je matice samých jedniček ($n\times n$).
Model \emph{náhodného surfaře}: s~pravděpodobností $d$ (typicky 0,85) následuje náhodný odkaz,
s~$1-d$ skočí na náhodnou stránku. \textbf{Proč augmentace:} řeší slepé uličky (stránky bez
odkazů) a~pasti --- rozšířená matice je stochastická, \emph{ireducibilní} i~\emph{aperiodická},
takže mocninná iterace konverguje k~jedinému stacionárnímu rozdělení. Druhý tvar vzniká
přeškálováním tak, aby $\mathbf e^{\!\top}P=n$.}

\podkap{16}{Detekce komunit}

\fml{Modularita $Q$}
\[
Q=\frac{1}{2m}\sum_{i,j}\Big(A_{ij}-\frac{k_ik_j}{2m}\Big)\delta(c_i,c_j)
\]
\pop{$A_{ij}$ = prvek matice sousednosti, $m$ = počet hran, $k_i$ = stupeň vrcholu $i$,
$c_i$ = komunita $i$, $\delta$ = Kroneckerovo delta (sčítá se tedy jen přes dvojice
\emph{uvnitř} komunit). $k_ik_j/2m$ je \emph{očekávaný} počet hran mezi $i$ a~$j$ v~náhodné
síti se stejnými stupni. Síť má smysluplnou komunitní strukturu, je-li hran uvnitř komunit
\emph{víc}, než by odpovídalo náhodě. Maximalizuje ji \emph{Louvainova metoda}
(hladově, $O(n\log n)$) a~její upřesnění \emph{Leiden}.}

\fml{Poměrový a~normalizovaný řez}
\[
\mathrm{RatioCut}(\pi)=\frac{1}{k}\sum_{i=1}^{k}\frac{\mathrm{cut}(C_i,\bar C_i)}{|C_i|},
\qquad
\mathrm{NCut}(\pi)=\frac{1}{k}\sum_{i=1}^{k}\frac{\mathrm{cut}(C_i,\bar C_i)}{\mathrm{vol}(C_i)}
\]
\pop{$\pi=\{C_1,\dots,C_k\}$ je \emph{rozdělení} na komunity (nikoli Ludolfovo číslo),
$\mathrm{cut}(A,B)=\sum_{i\in A, j\in B}a_{ij}$, $|C_i|$ = počet uzlů,
$\mathrm{vol}(C_i)$ = součet stupňů v~$C_i$. Prostý minimální řez vrací nevyvážená rozdělení
(často singleton) --- \textbf{obě míry proto normalizují velikostí komunity a~preferují
vyvážené řezy}. S~balančním omezením a~obecnými vahami je úloha \textbf{NP-těžká};
řeší se spektrálně.}

\podkap{17}{Predikce vazeb}

\fml{Sousedské míry predikce vazeb}
\[
\mathrm{CN}(i,j)=|S_i\cap S_j|,
\qquad
\mathrm{Jaccard}=\frac{|S_i\cap S_j|}{|S_i\cup S_j|},
\qquad
\mathrm{AA}(i,j)=\sum_{k\in S_i\cap S_j}\frac{1}{\log|S_k|}
\]
\pop{$S_i$ = sousedé uzlu $i$. Každá míra opravuje slabinu předchozí: \emph{společní sousedé}
nezohledňují stupně (dvě celebrity mají mnoho společných sousedů náhodou); \emph{Jaccard}
normalizuje vůči stupňům \emph{koncových} uzlů, ne však prostředních; \emph{Adamic--Adar}
snižuje váhu společného souseda s~jeho stupněm --- populární soused je společný mnoha párům,
a~je tedy méně informativní.}

\fml{Katzova míra}
\[
\mathrm{Katz}(i,j)=\sum_{t=1}^{\infty}\beta^{t}\,W_{ij}^{(t)},
\qquad
K=\sum_{t=1}^{\infty}\beta^{t}A^{t}=(I-\beta A)^{-1}-I
\]
\pop{$W_{ij}^{(t)}$ = počet procházek délky $t$ mezi $i$ a~$j$, $\beta<1$ = \emph{diskontní
faktor} potlačující dlouhé procházky (nutný pro konvergenci geometrické řady).
Měří \emph{nepřímou} propojenost, a~pomáhá proto ve \textbf{velmi řídkých} sítích, kde je
společných sousedů málo a~sousedské míry nerozlišují. $A$ je matice sousednosti, $I$ jednotková matice a~$K$ matice Katzových koeficientů. Součet Katzových koeficientů uzlu vůči
ostatním je jeho \emph{Katzova centralita}.}

\bigskip
\begin{center}
\begin{minipage}{0.9\textwidth}
\begin{poznbox}[Na co si dát pozor u~vzorců]
\small
\begin{itemize}
\item \textbf{Konvence pořadí:} u~turnajové selekce značí $i=1$ nejlepšího, u~lineární
  pořadové selekce nejhoršího. Zkoušející se na to rád zeptá.
\item \textbf{Směr optimalizace:} entropie $H(A)$ a~Gini se \emph{minimalizují}, informační
  zisk, $\chi^2$, MI a~$\Phi$ se \emph{maximalizují}.
\item \textbf{Off-policy vs.\ on-policy:} jediný rozdíl mezi Q-learningem a~SARSA je
  $\max_{a'}Q(s',a')$ vs.\ $Q(s',a')$.
\item \textbf{Harmonický vs.\ aritmetický průměr:} $F_1$ je harmonický, a~proto je blízko
  té \emph{menší} z~obou měr.
\item \textbf{Kde je normalizace nutná:} HITS, vlastní centralita a~PageRank bez ní
  nekonvergují; míry vzdálenosti bez normalizace atributů měří jen ten s~největším
  rozsahem.
\item \textbf{Nerovnost ve větě o~schématech} je \uv{$\ge$}, ne rovnost --- schéma může
  být křížením i~\emph{obnoveno}.
\end{itemize}
\end{poznbox}
\end{minipage}
\end{center}

\end{document}
